\chapter{Генеративное проектирование}

Предыдущая глава определила интерфейс между обучаемой моделью и физической задачей. Параметры конструкции обозначались через $\theta\in\Theta$, обучаемые параметры --- через $w$ или $\phi$, физическое состояние --- через $U_h$, проектный критерий --- через $J_h$, а условие задачи --- через $c\in\mathcal C$. Эти обозначения сохраняются. Новым объектом становится не ещё одна нейросеть, а механизм, который по условию задачи порождает множество кандидатов и позволяет исследовать пространство конструкций.

Генеративное проектирование в этой главе не отождествляется ни с топологической оптимизацией, ни с генеративной моделью. Топологическая оптимизация решает конкретную задачу поиска конструкции. Генеративная модель задаёт способ порождения кандидатов из распределения или латентного пространства. Генеративное проектирование объединяет постановку требований, представление конструкции, генерацию, физическую оценку, оптимизацию и проверку результата в один вычислительный процесс.

Сквозной контур главы имеет вид
\begin{equation}
\label{eq:generative-design-pipeline}
 (c,z)
 \xmapsto{\ G_\phi\ }
 \theta
 \xmapsto{\ \mathsf E\ }
 x_h
 \xmapsto{\ \mathcal H_{w,h}\ }
 \widehat U_h
 \longmapsto
 \widehat J_h
 \longmapsto
 \text{отбор и уточнение}
 \xmapsto{\ S_{h/2}\ }
 (U_{h/2},J_{h/2}).
\end{equation}
Здесь $z\in\mathcal Z$ --- латентный код или случайная переменная, $G_\phi$ --- генератор, $\mathsf E$ --- кодировщик конструкции для вычислений, $\mathcal H_{w,h}$ --- быстрый физический или гибридный оператор, а $S_{h/2}$ --- независимый расчёт на проверочной дискретизации. Если генератор или суррогат отсутствует, соответствующий блок просто заменяется точным отображением. Поэтому формула~\eqref{eq:generative-design-pipeline} описывает общий интерфейс, а не одну обязательную архитектуру.

Для условия $c$ множество допустимых конструкций обозначается
\begin{equation}
 \Theta_{\mathrm{adm}}(c)
 =
 \left\{
  \theta\in\Theta:
  g_i(\theta,U_h(\theta);c)\leq0,
  \quad
  h_j(\theta,U_h(\theta);c)=0
 \right\}.
\end{equation}
Функции $g_i$ задают неравенства, $h_j$ --- равенства. Они могут выражать ограничение объёма материала, минимальную толщину, допустимое напряжение, геометрическую совместимость или технологическое условие. Требование считается выполненным только после указания численного критерия проверки; наличие соответствующего штрафа в функции обучения само по себе допустимость не доказывает.

\section{Задача генеративного проектирования и пространство допустимых конструкций}

Генеративное проектирование начинается с обычной инженерной постановки. До появления генератора необходимо определить, что считается конструкцией, какие физические уравнения задают её отклик, какой критерий требуется уменьшать и какие проекты запрещены. Если эти четыре части не зафиксированы, генеративная модель не имеет математически определённой цели: она может выдавать разнообразные объекты, но нельзя проверить, решают ли они исходную задачу.

Для условия задачи $c\in\mathcal C$ конструкция задаётся параметрами $\theta\in\Theta$. Условие $c$ содержит то, что не выбирается алгоритмом проектирования: нагрузки, закрепления, допустимый объём материала, характеристики среды и другие требования конкретного расчёта. Параметры $\theta$ содержат то, что разрешено менять. Физическое состояние определяется уравнением
\begin{equation}
\label{eq:design-state-equation}
 R_h(U_h,\theta;c)=0,
\end{equation}
а проектный критерий вычисляется как
\begin{equation}
 J_h=J_h(\theta,U_h;c).
\end{equation}
Таким образом, зависимость критерия от конструкции проходит через физическое состояние:
\begin{equation}
\label{eq:reduced-design-objective}
 \widehat J_h(\theta;c)
 =
 J_h\bigl(\theta,U_h(\theta;c);c\bigr).
\end{equation}
Именно функция~\eqref{eq:reduced-design-objective}, а не внешний вид конструкции, является объектом инженерной оптимизации.

\subsection{Проектные переменные, условие и физическое состояние}

Разделение $\theta$ и $c$ принципиально для последующей условной генерации. Рассмотрим линейно-упругую конструкцию. Тогда дискретное состояние $U_h$ находится из уже введённой в главах~\ref{chap:partial-differential-equations} и~\ref{chap:continuum-mechanics} системы
\begin{equation}
\label{eq:elastic-state-design}
 K(\theta;c)U_h=F(c).
\end{equation}
Матрица $K$ зависит от геометрии и распределения материала, поэтому зависит от $\theta$. Вектор $F$ определяется нагрузкой и потому относится к условию $c$, если нагрузка не является проектной переменной.

Одна и та же физическая конструкция может участвовать в разных задачах. При фиксированном $\theta$ изменение $c$ меняет, например, точку приложения силы или набор закреплённых степеней свободы и тем самым меняет $U_h$. Наоборот, при фиксированном $c$ изменение $\theta$ меняет сам проект. В генеративной модели эта разница станет контрактом
\begin{equation}
 (z,c)\xmapsto{\ G_\phi }\theta,
\end{equation}
то есть условие передаётся генератору, а конструкция является его выходом.

Параметр не становится проектной переменной только потому, что его можно записать числом. Если материал заранее задан техническим заданием, модуль Юнга входит в $c$. Если алгоритму разрешён выбор материала, тот же параметр или дискретный индекс материала становится частью $\theta$. Граница между $c$ и $\theta$ задаётся постановкой задачи, а не типом математического объекта.

\subsection{Целевой функционал и податливость}

В статической линейной упругости базовым критерием будет податливость
\begin{equation}
 J_{\mathrm{comp}}(\theta;c)
 =F(c)^{\trans}U_h(\theta;c).
\end{equation}
При выполнении~\eqref{eq:elastic-state-design} и симметричной матрице жёсткости
\begin{equation}
 J_{\mathrm{comp}}
 =U_h^{\trans}K(\theta;c)U_h.
\end{equation}
Чем меньше $J_{\mathrm{comp}}$ при одной и той же нагрузке, тем меньше перемещение в направлении нагрузки и тем жёстче конструкция в рассматриваемом режиме. Поэтому задача минимизации податливости является удобной проверяемой базой для всей главы.

Именно такая постановка используется у Sigmund: проект описывается относительными плотностями конечных элементов, равновесие записывается как $KU=F$, а целевая функция как $U^{\trans}KU$. В его обозначениях одновременно задаются ограничение объёма и границы плотностей. Мы сохраняем эту структуру, но используем сквозные обозначения проекта: $J_{\mathrm{comp}}$ вместо $c(x)$, чтобы буква $c$ оставалась условием генеративной задачи.

Податливость не является универсальной мерой качества конструкции. Ограничения по напряжению, устойчивости, собственным частотам, массе и технологичности требуют других функционалов или дополнительных ограничений. В этой главе $J_{\mathrm{comp}}$ используется как минимальный механический пример, для которого физический решатель, значение критерия и производная по конструкции могут быть прослежены до конкретных формул.

\subsection{Ограничения и допустимое множество}

Без ограничений задача проектирования обычно вырождается. Если требуется только минимизировать податливость и разрешено бесконечно увеличивать количество материала, оптимум достигается добавлением материала, а не поиском формы. Поэтому вместе с критерием задаётся множество допустимых проектов
\begin{equation}
 \Theta_{\mathrm{adm}}(c)
 =
 \left\{
 \theta\in\Theta:
 g_i(\theta,U_h;c)\leq0,
 \quad
 h_j(\theta,U_h;c)=0
 \right\}.
\end{equation}
Здесь $g_i$ и $h_j$ проверяются на физическом состоянии, найденном из~\eqref{eq:design-state-equation}. Следовательно, ограничение может быть геометрическим, материальным или физическим.

Для плотностной задачи Sigmund ограничивает долю материала. Если $V_0$ --- объём всей проектной области, а $V(\rho)$ --- объём материала при плотностях $\rho$, то
\begin{equation}
\label{eq:volume-fraction}
 \frac{V(\rho)}{V_0}=f_{\mathrm{vol}},
 \qquad
 0<f_{\mathrm{vol}}<1.
\end{equation}
При одинаковых объёмах конечных элементов это условие становится ограничением на среднюю плотность:
\begin{equation}
\label{eq:mean-density}
 \frac{1}{N}\sum_{e=1}^{N}\rho_e=f_{\mathrm{vol}}.
\end{equation}
В вычислениях дополнительно используется нижняя граница
\begin{equation}
\label{eq:density-bounds}
 0<\rho_{\min}\leq\rho_e\leq1,
\end{equation}
чтобы глобальная матрица жёсткости не становилась сингулярной из-за элементов с ровно нулевой жёсткостью.

Условие допустимости нельзя заменить штрафом в функции обучения. Если генератор обучался с членом, который наказывает превышение объёма, итоговая конструкция всё равно должна пройти явную проверку~\eqref{eq:volume-fraction}. Это различие будет сохраняться во всей главе: loss управляет обучением, а $\Theta_{\mathrm{adm}}(c)$ определяет инженерную допустимость результата.

\subsection{Параметрическая оптимизация, форма и топология}

Слова ``параметр'', ``форма'' и ``топология'' описывают разные уровни свободы проектирования. Пусть конструкция занимает область $\Omega(\theta)\subset D$.

В параметрической оптимизации заранее выбирается конечномерное семейство
\begin{equation}
 \theta\in\R^m
 \longmapsto
 \Omega(\theta).
\end{equation}
Например, $\theta$ может содержать толщину стенки и радиус отверстия. Алгоритм ищет лучший объект только внутри заранее заданного семейства. Если параметризация всегда описывает одну пластину с одним отверстием, оптимизация не сможет породить второе отверстие или удалить первое.

При оптимизации формы меняется граница $\partial\Omega$, но обычно сохраняется тип связности области. Математически проектная переменная может описывать смещение границы или поле, задающее её положение. Такой класс шире простой параметрической геометрии, однако изменение числа компонент или появление нового отверстия уже требует дополнительного механизма.

В топологической оптимизации разрешается менять распределение материала внутри всей проектной области $D$. Плотностная запись
\begin{equation}
 \rho_h:D\to[\rho_{\min},1]
\end{equation}
после дискретизации даёт по одной проектной переменной на элемент. Области с $\rho_e\approx1$ интерпретируются как материал, а области с $\rho_e\approx\rho_{\min}$ --- как пустота. Поэтому число отверстий, ветвей и связных компонент не фиксируется небольшим вектором геометрических размеров.

Это различие важно для генеративного проектирования. Генератор не может исследовать те изменения конструкции, которых нет в его выходном представлении. Если $G_\phi$ выдаёт три геометрических размера, он остаётся генератором параметрических проектов. Если выходом является поле плотности или неявное поле, пространство потенциальных топологий существенно шире. Подробно представления сравниваются в разделе~\ref{sec:design-space-representation}.

Идеальная топологическая переменная бинарна: материал либо есть, либо отсутствует. Для элемента это означало бы $\rho_e\in\{0,1\}$, что приводит к дискретной комбинаторной задаче. SIMP заменяет её непрерывной релаксацией~\eqref{eq:density-bounds} и задаёт элементную жёсткость степенным законом
\begin{equation}
\label{eq:simp-element-stiffness}
 K_e(\rho_e)=\rho_e^p K_e^0,
 \qquad p>1.
\end{equation}
Здесь $K_e^0$ --- матрица жёсткости элемента из сплошного материала. В исходной реализации Sigmund типичным значением является $p=3$.

Степень $p>1$ делает промежуточную плотность относительно невыгодной по жёсткости. Например, при $p=3$ элемент с $\rho_e=1/2$ использует половину полного объёма элемента, но сохраняет только
\begin{equation}
 \left(\frac12\right)^3=\frac18
\end{equation}
его полной жёсткости. Поэтому оптимизация стремится перераспределять материал к значениям, близким к двум концам допустимого интервала.

Непрерывная релаксация сама по себе не задаёт масштаб деталей. Sigmund дополняет её фильтрацией чувствительностей, чтобы подавить сеточно-зависимые мелкие структуры. Поэтому SIMP определяется не одной интерполяцией~\eqref{eq:simp-element-stiffness}: физическая модель, объёмное ограничение, штраф промежуточных плотностей и регуляризация вместе задают задачу, результат которой затем можно использовать как обучающий пример или как старт для дальнейшего уточнения.

\subsection{Приведённая задача оптимизации}

После решения уравнения состояния $U_h=U_h(\theta;c)$ ограниченная физическая задача записывается только через проектные переменные:
\begin{equation}
\label{eq:reduced-constrained-problem}
 \theta^\star(c)
 =
 \argmin_{\theta\in\Theta_{\mathrm{adm}}(c)}
 \widehat J_h(\theta;c).
\end{equation}
Эта запись скрывает решение физической системы внутри функции $\widehat J_h$, но не устраняет его вычислительную стоимость. При каждом новом $\theta$ требуется получить состояние, если только точный решатель не заменён проверенным суррогатом.

Для SIMP-задачи минимизации податливости получаем конкретный вариант
\begin{equation}
 \begin{aligned}
 \min_{\rho}\quad &J_{\mathrm{comp}}(\rho)=F^{\trans}U,\\
 \text{при}\quad &K(\rho)U=F,\\
 &\frac{1}{N}\sum_{e=1}^{N}\rho_e=f_{\mathrm{vol}},\\
 &\rho_{\min}\leq\rho_e\leq1.
 \end{aligned}
\end{equation}
Это та же структура, что в постановке Sigmund: минимизация compliance, равновесие, заданная доля материала и ограниченные элементные плотности. Для главы 9 она важна не как конечный алгоритм, а как эталонная инженерная задача, относительно которой можно определить смысл генерации.

\subsection{Законченный расчёт: два параллельных элемента}

Рассмотрим минимальную систему, в которой penalization уже меняет проектное решение. Один узел с перемещением $u$ удерживается двумя параллельными упругими элементами. Полная жёсткость каждого элемента при $\rho_e=1$ равна $k_0=1$, внешняя сила равна $F=1$, а SIMP-показатель $p=3$. Плотности $\rho_1$ и $\rho_2$ являются проектными переменными.

Глобальная жёсткость системы равна сумме двух параллельных ветвей:
\begin{equation}
 K(\rho_1,\rho_2)
 =\rho_1^3+\rho_2^3.
\end{equation}
Из равновесия $Ku=F$ получаем
\begin{equation}
 u
 =
 \frac{1}{\rho_1^3+\rho_2^3}.
\end{equation}
Так как $F=1$, податливость совпадает с перемещением:
\begin{equation}
 J_{\mathrm{comp}}
 =Fu
 =
 \frac{1}{\rho_1^3+\rho_2^3}.
\end{equation}

Пусть разрешено использовать половину полного объёма двух одинаковых элементов. Тогда $f_{\mathrm{vol}}=1/2$, и из~\eqref{eq:mean-density}
\begin{equation}
\label{eq:two-element-volume}
 \rho_1+\rho_2=1.
\end{equation}
Сначала распределим материал равномерно:
\begin{equation}
 \rho_1=\rho_2=0.5.
\end{equation}
Тогда
\begin{equation}
 K_{\mathrm{uniform}}
 =2(0.5)^3
 =0.25,
 \qquad
 J_{\mathrm{uniform}}=4.
\end{equation}

Теперь почти весь разрешённый материал перенесём в первую ветвь. При $\rho_{\min}=10^{-3}$ можно взять
\begin{equation}
 \rho_1=0.999,
 \qquad
 \rho_2=0.001,
\end{equation}
что сохраняет то же ограничение~\eqref{eq:two-element-volume}. Получаем
\begin{equation}
 \begin{aligned}
 K_{\mathrm{conc}}
 &=0.999^3+0.001^3
 \approx0.997003,\\
 J_{\mathrm{conc}}
 &=\frac{1}{K_{\mathrm{conc}}}
 \approx1.003006.
 \end{aligned}
\end{equation}
При одинаковом количестве материала податливость уменьшилась примерно в
\begin{equation}
 \frac{J_{\mathrm{uniform}}}{J_{\mathrm{conc}}}
 \approx3.99
\end{equation}
раза. Причина видна непосредственно из кубического закона: две половинные плотности дают суммарную жёсткость $1/4$, тогда как одна почти сплошная ветвь даёт жёсткость почти $1$.

Расчёт намеренно минимален: одна степень свободы состояния и два элемента. Он отделяет четыре слоя постановки. Плотности $\rho$ задают проект; уравнение $Ku=F$ --- физику; $J_{\mathrm{comp}}$ --- цель; условие $\rho_1+\rho_2=1$ --- ограничение. На конечномерной сетке тот же принцип дополняется чувствительностью, фильтром и итерационным обновлением плотностей.

Задача~\eqref{eq:reduced-constrained-problem} отвечает на вопрос: какой проект минимизирует выбранный критерий при фиксированном условии $c$ и выбранной процедуре поиска. Генеративное проектирование ставит дополнительную задачу: быстро получать несколько различных кандидатов для разных $c$, сохраняя инженерную допустимость и качество.

Поэтому генератор не заменяет определение оптимума. Он вводит распределение предложений
\begin{equation}
 z\sim p_Z,
 \qquad
 \theta^{(m)}=G_\phi(z_m,c),
 \qquad m=1,\ldots,M,
\end{equation}
которые затем оцениваются относительно той же физической постановки. Для каждого кандидата остаются обязательными
\begin{equation}
 R_h(U_h^{(m)},\theta^{(m)};c)=0,
 \qquad
 \theta^{(m)}\in\Theta_{\mathrm{adm}}(c),
 \qquad
 J_h^{(m)}=J_h(\theta^{(m)},U_h^{(m)};c).
\end{equation}
Если эти величины не вычислены или не проверены, выход $G_\phi$ остаётся геометрическим предложением, а не подтверждённой инженерной конструкцией.

Для задач типа рассматриваемых в AIRI это и задаёт полезную границу между машинным обучением и физическим проектированием. ML-блок может сокращать поиск, кодировать сложное пространство геометрий или давать несколько стартовых точек. Уравнения состояния, ограничения и независимая физическая проверка определяют, какие из этих предложений действительно решают задачу. Следующий раздел уточняет, какие представления $\theta$ позволяют генератору изменять размеры, форму и топологию, не смешивая эти три режима.

\section{Представление пространства конструкций}
\label{sec:design-space-representation}

Генеративная модель не изменяет физический объект непосредственно. Она изменяет координаты выбранного представления, после чего эти координаты преобразуются в геометрию, поле материала или расчётную сетку. Поэтому представление конструкции является частью постановки задачи: оно определяет, какие проекты вообще достижимы, какие ограничения удобно проверять и через какие величины можно передавать градиент.

Пусть $\Theta$ --- пространство записей конструкции, а $\mathcal G$ --- пространство геометрических объектов или распределений материала. Введём отображение интерпретации
\begin{equation}
\label{eq:representation-to-geometry}
 \mathsf G:\Theta\longrightarrow\mathcal G,
 \qquad
 \theta\longmapsto \mathsf G(\theta).
\end{equation}
Физический решатель обычно не работает непосредственно с $\mathsf G(\theta)$. Ему требуется дискретное представление $x_h$, поэтому используется ещё одно отображение
\begin{equation}
 \theta
 \xmapsto{\ \mathsf E_h\ }
 x_h
 \longmapsto
 K(x_h),F(c)
 \longmapsto
 U_h.
\end{equation}
Здесь $\mathsf E_h$ включает операции, необходимые для расчёта: построение сетки, назначение материалу конечных элементов, вычисление геометрических коэффициентов или сборку другой дискретной модели. В главе~\ref{chap:connections-machine-learning} уже обсуждалось представление геометрии для обучаемой модели. Здесь тот же вопрос рассматривается со стороны пространства поиска: какие конструкции остаются внутри образа выбранной параметризации.

Если алгоритм меняет $\theta$, то фактически он ищет не во всём множестве мыслимых конструкций, а в образе отображения~\eqref{eq:representation-to-geometry}:
\begin{equation}
\label{eq:reachable-design-set}
 \mathcal G_{\mathrm{reach}}
 =
 \mathsf G(\Theta)
 \subseteq
 \mathcal G.
\end{equation}
Ограничение~\eqref{eq:reachable-design-set} существует до выбора оптимизатора. Градиентный спуск, эволюционный поиск и генеративная модель могут исследовать разные части $\Theta$, но ни один из них не создаст геометрию вне $\mathsf G(\Theta)$.

Поэтому уменьшение размерности имеет две стороны. Небольшое число координат удешевляет поиск и позволяет встроить известные геометрические ограничения прямо в параметризацию. Одновременно оно исключает часть конструкций. Для генеративного проектирования это особенно важно: разнообразие выборки ограничено не только распределением $p_\phi$, но и самим множеством $\mathcal G_{\mathrm{reach}}$.

Работа Oh et al. рассматривает design representation и shape parameterization как отдельную часть глубокого генеративного проектирования и приводит примеры низкоразмерных представлений через автоэнкодеры. Для нас существенен не конкретный тип сети, а следствие: переход к скрытым координатам меняет множество доступных проектов и потому должен проверяться так же явно, как физические ограничения.

\subsection{Параметрическая геометрия}

В конечномерной параметризации конструкция задаётся вектором
\begin{equation}
 \theta=(\theta_1,\ldots,\theta_m)\in\Theta\subset\R^m,
 \qquad
 \Omega=\Omega(\theta).
\end{equation}
Компоненты могут задавать длины, толщины, радиусы, углы или коэффициенты базисного разложения границы. Если, например, пластина описывается шириной $b$, высотой $h$ и радиусом одного отверстия $r$, то
\begin{equation}
 \theta=(b,h,r)
\end{equation}
задаёт трёхмерное пространство поиска независимо от числа узлов расчётной сетки.

Преимущество такого представления состоит в явном смысле каждой координаты. Ограничение минимальной толщины можно записать непосредственно как неравенство по $\theta$, а гладкая зависимость $\Omega(\theta)$ позволяет вычислять производные проектного критерия по параметрам. Недостаток также структурный: число отверстий, ветвей и компонент связности обычно фиксировано формой параметризации.

Если физический критерий записан как $\widehat J_h(\theta;c)$, то локальное изменение проекта имеет вид
\begin{equation}
\label{eq:parametric-first-variation}
 \widehat J_h(\theta+\delta\theta;c)
 =
 \widehat J_h(\theta;c)
 +
 \nabla_\theta \widehat J_h(\theta;c)^{\trans}\delta\theta
 +o(\|\delta\theta\|).
\end{equation}
Формула~\eqref{eq:parametric-first-variation} полезна только внутри выбранного семейства. Нулевой градиент означает локальную стационарность по доступным параметрам, но не оптимальность среди геометрий, которых параметризация не умеет представлять.

\subsection{Плотностное представление на фиксированной сетке}

В SIMP проектная область $D$ заранее разбивается на $N$ конечных элементов. Конструкция задаётся вектором плотностей
\begin{equation}
 \rho=(\rho_1,\ldots,\rho_N)
 \in
 [\rho_{\min},1]^N.
\end{equation}
Ему соответствует кусочно-постоянное поле
\begin{equation}
 \rho_h(x)
 =
 \sum_{e=1}^{N}\rho_e\chi_e(x),
\end{equation}
где $\chi_e$ равна единице на элементе $e$ и нулю вне него. В работе Sigmund именно относительные плотности элементов являются проектными переменными; через степенной закон они определяют элементные жёсткости и далее глобальную систему равновесия.

Для физического решателя отображение $\mathsf E_h$ здесь почти тождественно: сетка уже задана, а каждая координата $\rho_e$ относится к конкретному элементу. Глобальная матрица имеет структуру
\begin{equation}
 K(\rho)
 =
 \sum_{e=1}^{N}
 A_e^{\trans}
 \rho_e^p K_e^0
 A_e,
\end{equation}
где $A_e$ размещает локальные степени свободы элемента в глобальном векторе. Поэтому плотностный массив является не картинкой материала, а координатной записью коэффициентов физической дискретизации.

Размерность такого пространства растёт вместе с сеткой. При удвоении разрешения по каждой координате двумерной области число проектных переменных увеличивается примерно в четыре раза. Это даёт свободу изменения топологии, но одновременно делает необходимыми фильтрацию масштаба деталей, согласование разных сеток и аккуратную проверку того, что генератор не запомнил конкретную нумерацию элементов.

\subsection{Неявное поле и движущаяся граница}

Неявное представление задаёт скалярное поле
\begin{equation}
 \psi_\theta:D\to\R,
\end{equation}
по которому материал определяется знаком:
\begin{equation}
 \Omega(\theta)
 =
 \{x\in D:\psi_\theta(x)<0\},
 \qquad
 \partial\Omega(\theta)
 =
 \{x:\psi_\theta(x)=0\}.
\end{equation}
Такое представление отделяет геометрическую границу от явного списка её вершин. Изменение параметров поля перемещает нулевой уровень, а при подходящей семье $\psi_\theta$ компоненты нулевого множества могут появляться, сливаться и исчезать.

Это не означает автоматической дифференцируемости всего расчётного контура. Даже если $\psi_\theta(x)$ гладко зависит от $\theta$, операция построения сетки по нулевому уровню может содержать дискретные переключения. Поэтому необходимо различать производную поля
\begin{equation}
 \frac{\partial\psi_\theta(x)}{\partial\theta}
\end{equation}
и производную итогового критерия после геометрической дискретизации. Аналогичное различие уже появлялось в главе~\ref{chap:connections-machine-learning} при обсуждении неявных геометрий: гладкая запись границы ещё не гарантирует гладкость всего программного пути до $J_h$.

Для генеративной модели неявное поле удобно тем, что выход может быть определён на фиксированной области $D$, а геометрия извлекается после генерации. Но допустимость всё равно проверяется на геометрическом объекте: минимальная толщина, связность и отсутствие изолированных компонент не следуют из одного факта, что поле имеет корректный размер массива.

\subsection{Латентное представление и образ декодера}

Пусть высокоразмерная конструкция $\theta\in\Theta$ получается из скрытого кода $z\in\mathcal Z\subset\R^d$ через декодер
\begin{equation}
 D_\phi:\mathcal Z\longrightarrow\Theta,
 \qquad
 \theta=D_\phi(z),
 \qquad
 d\ll\dim\Theta.
\end{equation}
Тогда реальное пространство поиска равно не всей $\Theta$, а
\begin{equation}
\label{eq:decoder-image}
 \Theta_\phi
 =
 D_\phi(\mathcal Z)
 \subseteq
 \Theta.
\end{equation}
Oh et al. обсуждают использование VAE и других генеративных моделей для кодирования высокоразмерных проектных переменных в низкоразмерное пространство. С точки зрения нашей постановки это означает замену исходной задачи на задачу, ограниченную образом~\eqref{eq:decoder-image}.

Если оптимизация выполняется по $z$, приведённый критерий имеет вид
\begin{equation}
 \widetilde J_h(z;c)
 =
 \widehat J_h(D_\phi(z);c).
\end{equation}
Для дифференцируемого декодера правило цепочки даёт
\begin{equation}
 \nabla_z\widetilde J_h
 =
 \bigl(DD_\phi(z)\bigr)^{\trans}
 \nabla_\theta\widehat J_h.
\end{equation}
Матрица $DD_\phi(z)$ переводит изменение скрытого кода в допустимое для декодера изменение конструкции. Поэтому даже точный физический градиент $\nabla_\theta\widehat J_h$ проецируется на направления, существующие внутри обученного представления.

Низкая размерность полезна, если образ декодера содержит нужные семейства хороших конструкций. Если обучающие данные не содержали определённый тип геометрии, декодер может исключить его из $\Theta_\phi$. Улучшение оптимизатора в этом случае не устраняет ограничение представления.

Для главы 9 удобно явно разделить три объекта:
\begin{equation}
 \theta
 \xmapsto{\ \mathsf G\ }
 \Omega\ \text{или}\ \rho
 \xmapsto{\ \mathsf E_h\ }
 x_h
 \xmapsto{\ S_h\ }
 U_h.
\end{equation}
Первая стрелка задаёт смысл проектных координат. Вторая строит данные конкретной численной схемы. Третья решает физическую задачу. Такое разделение предотвращает типичную подмену, когда массив размера $n_x\times n_y$ одновременно называют и геометрией, и сеткой, и входом нейросети.

При фиксированной плотностной сетке первые две стрелки почти совпадают. Для CAD-параметров $\mathsf E_h$ должна построить сетку заново или деформировать существующую. Для неявного поля требуется извлечь границу или назначить материал ячейкам. Для латентного кода перед физическим интерфейсом появляется декодер. Следовательно, вычислительная стоимость и дифференцируемость проектного контура зависят не только от $S_h$, но и от $\mathsf G$ и $\mathsf E_h$.

Для независимой проверки полезно менять дискретизацию, не меняя физическую конструкцию. Если $x_h=\mathsf E_h(\theta)$ и $x_{h/2}=\mathsf E_{h/2}(\theta)$ описывают один и тот же проект, то сравнение
\begin{equation}
 J_h(\theta)
 \quad\text{и}\quad
 J_{h/2}(\theta)
\end{equation}
проверяет численную устойчивость результата. Если же сама $\theta$ привязана к пикселям конкретной сетки, перед таким сравнением необходимо определить перенос конструкции между разрешениями. Без этого изменение $h$ одновременно меняет и дискретизацию физики, и сам проект.

\subsection{Законченный расчёт: ограничение образа декодера}

Рассмотрим проект из четырёх элементных плотностей
\begin{equation}
 \rho=(\rho_1,\rho_2,\rho_3,\rho_4)\in[0,1]^4.
\end{equation}
Пусть целевая бинарная топология содержит материал в первом и четвёртом элементах:
\begin{equation}
\label{eq:latent-target}
 \rho^\star=(1,0,0,1).
\end{equation}
Она использует две единицы материала, то есть удовлетворяет условию
\begin{equation}
 \sum_{e=1}^{4}\rho_e=2.
\end{equation}

Предположим теперь, что обученное низкоразмерное представление имеет два параметра $a,b\in[0,1]$ и декодирует их как
\begin{equation}
\label{eq:tied-decoder}
 D(a,b)=(a,b,a,b).
\end{equation}
Такой декодер связывает плотности первого и третьего элементов, а также второго и четвёртого. Возможно, эта симметрия возникла из структуры обучающих данных. Требуется найти ближайший к~\eqref{eq:latent-target} проект, достижимый через~\eqref{eq:tied-decoder}.

Квадрат евклидовой ошибки равен
\begin{equation}
 \begin{aligned}
 E(a,b)
 &=\|D(a,b)-\rho^\star\|_2^2\\
 &=(a-1)^2+b^2+a^2+(b-1)^2\\
 &=2\left(a-\frac12\right)^2
   +2\left(b-\frac12\right)^2
   +1.
 \end{aligned}
\end{equation}
Минимум достигается при
\begin{equation}
 a=b=\frac12,
 \qquad
 D(a,b)=\left(\frac12,\frac12,\frac12,\frac12\right),
\end{equation}
и минимальная ошибка составляет
\begin{equation}
\label{eq:decoder-min-error}
 E_{\min}=1.
\end{equation}
Так как $\|\rho^\star\|_2^2=2$, относительная ошибка по норме равна
\begin{equation}
 \frac{\|D(a,b)-\rho^\star\|_2}
      {\|\rho^\star\|_2}
 =
 \frac{1}{\sqrt2}
 \approx0.707.
\end{equation}

При этом найденный проект использует ровно тот же объём материала:
\begin{equation}
 4\cdot\frac12=2.
\end{equation}
Следовательно, объёмное ограничение выполнено идеально, но нужная топология недостижима. Ни увеличение числа шагов оптимизации, ни точный градиент, ни более сильный штраф не устранят ошибку~\eqref{eq:decoder-min-error}, пока не изменён сам декодер. Этот расчёт показывает, почему качество представления необходимо отделять от качества оптимизатора и генеративной модели.

\subsection{Выбор представления для генеративного проектирования}

Практический выбор представления определяется тем, какие изменения конструкции действительно нужны. Если требуется подобрать несколько размеров известной детали, конечномерная параметризация даёт наиболее прозрачный контракт. Если требуется перераспределять материал и менять число отверстий или ветвей на фиксированной области, плотностное или неявное представление предоставляет существенно более широкое пространство. Если исходная размерность слишком велика, латентный код может уменьшить стоимость поиска, но вводит дополнительное ограничение $\Theta_\phi$.

Для исследовательского контура, связанного с машинным обучением и генеративным проектированием в задачах типа рассматриваемых в AIRI, важна не универсальная ``лучшая'' кодировка, а явный путь
\begin{equation}
 z,c
 \longrightarrow
 \theta
 \longrightarrow
 \mathsf G(\theta)
 \longrightarrow
 \mathsf E_h(\theta)
 \longrightarrow
 U_h
 \longrightarrow
 J_h.
\end{equation}
На каждой стрелке должно быть известно, какая информация теряется, какие ограничения сохраняются и через что проходит производная. Тогда генеративная модель можно менять независимо от физического решателя, а физическую проверку --- независимо от способа генерации.

Для плотностного представления детерминированной основой служит один полный шаг SIMP: от $\rho$ и сборки $K(\rho)$ до чувствительности, фильтра и обновлённого проекта. Тот же оператор затем можно применять к кандидатам генератора.

\section{Физическая оптимизация конструкции}
\label{sec:physics-based-design-optimization}

После выбора представления задача становится вычислительной: по текущей конструкции нужно получить физическое состояние, значение критерия и направление изменения проектных переменных. Для плотностной топологической оптимизации этот цикл имеет вид
\begin{equation}
 \rho
 \longmapsto
 K(\rho)
 \longmapsto
 U(\rho)
 \longmapsto
 J_{\mathrm{comp}}(\rho)
 \longmapsto
 \nabla_\rho J_{\mathrm{comp}}
 \longmapsto
 \rho^+.
\end{equation}
Здесь $\rho=(\rho_1,\ldots,\rho_N)$ --- плотности конечных элементов, а податливость обозначается $J_{\mathrm{comp}}$, чтобы не смешивать её с условием задачи $c$. Работа Sigmund удобна тем, что в ней весь этот контур записан в одной постановке: SIMP-интерполяция, уравнение равновесия, ограничение объёма, чувствительность, фильтр и правило обновления.

В этой секции не выводится заново метод конечных элементов. Глава~\ref{chap:continuum-mechanics} уже дала физический смысл матрицы жёсткости и вектора перемещений. Здесь требуется понять другой объект: как изменение материала в отдельном элементе проходит через физический решатель и превращается в изменение проектного критерия.

Пусть условие задачи $c$ фиксирует расчётную область, закрепления, нагрузку и допустимое количество материала. При заданной плотности $\rho$ дискретная задача линейной упругости имеет вид
\begin{equation}
\label{eq:simp-state-equation}
 K(\rho)U(\rho)=F(c).
\end{equation}
Вектор $U$ не является независимой проектной переменной: после выбора $\rho$ он определяется уравнением~\eqref{eq:simp-state-equation}. Поэтому исходный критерий $J_{\mathrm{comp}}(\rho,U;c)$ можно рассматривать как функцию только проектных переменных,
\begin{equation}
\label{eq:reduced-compliance}
 \widehat J_{\mathrm{comp}}(\rho;c)
 =
 J_{\mathrm{comp}}\bigl(\rho,U(\rho);c\bigr).
\end{equation}
Такой переход называется исключением состояния: физический решатель остаётся внутри функции~\eqref{eq:reduced-compliance}, хотя в обозначении $U$ больше не выписывается как отдельный аргумент.

Для статической линейной упругости податливость равна работе внешней нагрузки,
\begin{equation}
 J_{\mathrm{comp}}(\rho)
 =F^{\trans}U
 =U^{\trans}K(\rho)U.
\end{equation}
При фиксированном $F$ уменьшение податливости означает уменьшение перемещений в направлениях приложенных сил, то есть увеличение эффективной жёсткости конструкции относительно данного сценария нагружения. Это не универсальная мера ``качества детали'': изменение нагрузки меняет и физическую задачу, и смысл оптимума.

Практический цикл поэтому всегда начинается с нового решения~\eqref{eq:simp-state-equation}. Нельзя обновить плотности, оставить старый $U$ и считать полученное значение точной податливостью новой конструкции. После каждого изменения $\rho$ матрица $K(\rho)$ меняется, следовательно, состояние нужно пересчитать.

\subsection{SIMP-интерполяция жёсткости}

В плотностной релаксации каждому элементу назначается
\begin{equation}
\label{eq:simp-density-bounds}
 \rho_{\min}\leq \rho_e\leq 1,
\end{equation}
где малая положительная $\rho_{\min}$ предотвращает вырождение глобальной матрицы. Значение $\rho_e=1$ соответствует сплошному материалу, а значение около $\rho_{\min}$ --- почти пустому элементу.

В SIMP элементная матрица жёсткости задаётся степенным законом
\begin{equation}
 K_e(\rho_e)
 =
 \rho_e^p K_e^0,
 \qquad p>1,
\end{equation}
где $K_e^0$ --- матрица элемента из сплошного материала. После сборки
\begin{equation}
\label{eq:simp-global-stiffness}
 K(\rho)
 =
 \sum_{e=1}^{N}
 A_e^{\trans}\rho_e^pK_e^0A_e.
\end{equation}
При $p>1$ промежуточная плотность теряет жёсткость быстрее, чем материал. Например, при $p=3$ элемент с $\rho_e=1/2$ сохраняет половину условного объёма, но только
\begin{equation}
 \left(\frac12\right)^3=\frac18
\end{equation}
жёсткости сплошного элемента. Именно поэтому оптимизация получает стимул уходить от больших областей серого материала к значениям плотности, близким к концам допустимого интервала.

Степенной закон не является отдельной физической моделью нового материала, которую затем нужно интерпретировать буквально. В этой постановке он служит релаксацией дискретной задачи ``материал есть или материала нет''. Значение показателя $p$ влияет на геометрию ландшафта оптимизации, поэтому менять его без фиксации в эксперименте нельзя.

Если все элементы имеют одинаковый объём, заданная доля материала $f_{\mathrm{vol}}$ записывается как
\begin{equation}
 \frac1N\sum_{e=1}^{N}\rho_e
 \leq
 f_{\mathrm{vol}}.
\end{equation}
Для элементов разных объёмов $v_e$ используется
\begin{equation}
\label{eq:simp-volume-general}
 V(\rho)
 =
 \sum_{e=1}^{N}v_e\rho_e,
 \qquad
 V(\rho)\leq V_{\max}.
\end{equation}
Таким образом, оптимизация не должна просто добавлять материал туда, где он повышает жёсткость. Она должна перераспределять ограниченный ресурс.

Для активного ограничения~\eqref{eq:simp-volume-general} вводится множитель Лагранжа $\lambda$. Если временно отбросить границы~\eqref{eq:simp-density-bounds}, условие стационарности имеет вид
\begin{equation}
\label{eq:simp-kkt-stationarity}
 \frac{\partial \widehat J_{\mathrm{comp}}}{\partial\rho_e}
 +
 \lambda v_e
 =0.
\end{equation}
Формула~\eqref{eq:simp-kkt-stationarity} выражает баланс предельной выгоды. Если небольшая добавка материала в одном элементе уменьшает податливость значительно сильнее, чем в другом, текущая конструкция ещё не стационарна: материал выгодно переносить в первый элемент.

Множитель $\lambda$ не задаётся заранее. В алгоритме Sigmund он подбирается бисекцией так, чтобы обновлённые плотности удовлетворяли заданному объёму. Поэтому условие материала входит непосредственно в правило обновления, а не проверяется только после завершения оптимизации.

\subsection{Чувствительность податливости}

Нужно вычислить производную приведённого критерия. Для функции~\eqref{eq:reduced-compliance} начнём с
\begin{equation}
 J_{\mathrm{comp}}=F^{\trans}U.
\end{equation}
Нагрузка $F$ фиксирована, поэтому при изменении $\rho_e$
\begin{equation}
 \frac{\partial J_{\mathrm{comp}}}{\partial\rho_e}
 =
 F^{\trans}
 \frac{\partial U}{\partial\rho_e}.
\end{equation}
Дифференцируем уравнение равновесия~\eqref{eq:simp-state-equation}:
\begin{equation}
\label{eq:state-differentiation}
 \frac{\partial K}{\partial\rho_e}U
 +
 K\frac{\partial U}{\partial\rho_e}
 =0.
\end{equation}
Так как $F=KU$ и $K$ симметрична,
\begin{equation}
 \begin{aligned}
 \frac{\partial J_{\mathrm{comp}}}{\partial\rho_e}
 &=
 U^{\trans}K
 \frac{\partial U}{\partial\rho_e}\\
 &=-
 U^{\trans}
 \frac{\partial K}{\partial\rho_e}
 U.
 \end{aligned}
\end{equation}
Из~\eqref{eq:simp-global-stiffness}
\begin{equation}
 \frac{\partial K}{\partial\rho_e}
 =
 A_e^{\trans}
 p\rho_e^{p-1}K_e^0
 A_e.
\end{equation}
Обозначая локальный вектор перемещений через $u_e=A_eU$, получаем формулу Sigmund
\begin{equation}
\label{eq:simp-compliance-sensitivity}
 \boxed{
 \frac{\partial \widehat J_{\mathrm{comp}}}{\partial\rho_e}
 =
 -p\rho_e^{p-1}
 u_e^{\trans}K_e^0u_e
 }.
\end{equation}
Квадратичная форма $u_e^{\trans}K_e^0u_e$ неотрицательна, поэтому производная~\eqref{eq:simp-compliance-sensitivity} неположительна. При фиксированном состоянии добавление материала в отдельный элемент не увеличивает податливость. Чем больше локальная энергия деформации, тем сильнее алгоритм стремится сохранить или добавить материал в этой области.

Важно, что~\eqref{eq:simp-compliance-sensitivity} уже учитывает зависимость $U(\rho)$. Если дифференцировать только явный множитель $\rho_e^p$ в записи $U^{\trans}K(\rho)U$ и ошибочно считать $U$ постоянным, знак производной получится противоположным. Именно устранение производной состояния через~\eqref{eq:state-differentiation} делает формулу корректной.

\subsection{Фильтрация чувствительности и масштаб детали}

На фиксированной сетке соседние элементы можно переключать независимо. Без дополнительного масштаба это приводит к сеточно-зависимым мелким структурам и шахматным паттернам. В реализации Sigmund используется локальное усреднение чувствительностей.

Для элементов $e$ и $f$ введём вес
\begin{equation}
 H_{ef}
 =
 \max\bigl(0,r_{\min}-\operatorname{dist}(e,f)\bigr),
\end{equation}
где $r_{\min}$ задаёт радиус взаимодействия в единицах сетки. Отфильтрованная чувствительность записывается как
\begin{equation}
 \widetilde{\frac{\partial J}{\partial\rho_e}}
 =
 \frac{
  \displaystyle
  \sum_f H_{ef}\rho_f
  \frac{\partial J}{\partial\rho_f}
 }{
  \displaystyle
  \rho_e\sum_f H_{ef}
 }.
\end{equation}
В числителе участвуют только элементы внутри радиуса $r_{\min}$. Поэтому решение в одном элементе зависит от локального окружения, а не от единственной пиксельной координаты.

Фильтр меняет сам алгоритм поиска. Если изменить сетку, но оставить $r_{\min}$ численно тем же в количестве элементов, физический радиус фильтра изменится. Для воспроизводимого сравнения разрешений нужно фиксировать масштаб в геометрических единицах или явно пересчитывать его при изменении $h$.

Для генеративного проектирования это имеет дополнительное следствие. Генератор может выдать очень тонкие детали, которые выглядят допустимыми в массиве $\rho$, но затем исчезнут после фильтра или окажутся неустойчивыми при уточнении сетки. Поэтому размер минимальной детали относится к контракту представления и проверки, а не только к постобработке изображения.

\subsection{Правило обновления по критерию оптимальности}

После вычисления чувствительностей нужно получить новую плотность. В варианте optimality criteria, используемом Sigmund, вводится величина
\begin{equation}
\label{eq:simp-oc-ratio}
 B_e
 =
 -
 \frac{
  \widetilde{\partial J/\partial\rho_e}
 }{
  \lambda\,\partial V/\partial\rho_e
 }.
\end{equation}
При равных объёмах элементов знаменатель отличается между элементами только общей величиной $\lambda$. Чем больше $B_e$, тем выгоднее сохранить материал в элементе $e$.

До учёта границ естественное обновление имеет вид
\begin{equation}
\label{eq:simp-oc-unclipped}
 \rho_e^+
 =
 \rho_e B_e^{\eta},
 \qquad
 \eta=\frac12.
\end{equation}
На практике изменение ограничивается move-limit $m$:
\begin{equation}
 \rho_e^+
 =
 \operatorname{clip}
 \left(
  \rho_e B_e^{\eta},
  \max(\rho_{\min},\rho_e-m),
  \min(1,\rho_e+m)
 \right).
\end{equation}
Здесь $\operatorname{clip}(x,a,b)$ означает ограничение числа $x$ интервалом $[a,b]$. После пробного обновления проверяется объём. Если материала получилось слишком много, $\lambda$ увеличивают; если слишком мало --- уменьшают. Бисекция продолжается до выполнения объёмного ограничения с заданной точностью.

Move-limit нужен не для физического смысла плотности, а для устойчивости итерационного процесса: он не позволяет одной линейной аппроксимации чувствительности сразу перестроить конструкцию слишком сильно. После обновления снова решается система~\eqref{eq:simp-state-equation}, вычисляются новые энергии и начинается следующая итерация.

\subsection{Законченный расчёт: один шаг SIMP}

Рассмотрим три элемента одинакового объёма. На текущей итерации
\begin{equation}
 \rho^{(k)}
 =
 \left(\frac12,\frac12,\frac12\right),
 \qquad
 p=3,
\end{equation}
а после решения задачи упругости локальные квадратичные формы равны
\begin{equation}
\label{eq:simp-example-energies}
 u_e^{\trans}K_e^0u_e
 =
 \left(8,2,\frac12\right).
\end{equation}
По формуле~\eqref{eq:simp-compliance-sensitivity}
\begin{equation}
 \begin{aligned}
 \frac{\partial J}{\partial\rho_1}
 &=-3\left(\frac12\right)^2 8=-6,\\
 \frac{\partial J}{\partial\rho_2}
 &=-3\left(\frac12\right)^2 2=-\frac32,\\
 \frac{\partial J}{\partial\rho_3}
 &=-3\left(\frac12\right)^2\frac12=-\frac38.
 \end{aligned}
\end{equation}
Первый элемент наиболее чувствителен: единица добавленного туда материала обещает наибольшее локальное уменьшение податливости.

Чтобы увидеть само правило OC без дополнительных эффектов, в этом расчёте не применяем фильтр, move-limit и границы. Требуем сохранить исходный объём,
\begin{equation}
 \rho_1^++\rho_2^++\rho_3^+
 =
 \frac32.
\end{equation}
При равных объёмах элементов из~\eqref{eq:simp-oc-ratio}--\eqref{eq:simp-oc-unclipped}
\begin{equation}
\label{eq:simp-example-update}
 \rho_e^+
 =
 \frac12
 \sqrt{\frac{s_e}{\lambda}},
 \qquad
 s=\left(6,\frac32,\frac38\right).
\end{equation}
Подставляем~\eqref{eq:simp-example-update} в ограничение объёма:
\begin{equation}
 \frac{1}{2\sqrt\lambda}
 \left(
  \sqrt6+\sqrt{\frac32}+\sqrt{\frac38}
 \right)
 =
 \frac32.
\end{equation}
Так как
\begin{equation}
 \sqrt{\frac32}=\frac{\sqrt6}{2},
 \qquad
 \sqrt{\frac38}=\frac{\sqrt6}{4},
\end{equation}
получаем
\begin{equation}
 \sqrt\lambda
 =
 \frac{7\sqrt6}{12},
 \qquad
 \lambda=\frac{49}{24}.
\end{equation}
Новые плотности равны
\begin{equation}
 \boxed{
 \rho^{(k+1)}
 =
 \left(
  \frac67,
  \frac37,
  \frac3{14}
 \right)
 }.
\end{equation}
Их сумма действительно равна $3/2$. Материал перенесён из третьего и частично второго элемента в первый, где чувствительность была наибольшей.

Линейная оценка изменения критерия составляет
\begin{equation}
 \begin{aligned}
 \Delta J
 &\approx
 \nabla_\rho J^{\trans}
 \bigl(\rho^{(k+1)}-\rho^{(k)}\bigr)\\
 &=-6\frac5{14}
   -\frac32\left(-\frac1{14}\right)
   -\frac38\left(-\frac4{14}\right)\\
 &=-\frac{27}{14}<0.
 \end{aligned}
\end{equation}
Это только локальный прогноз. Точное $J_{\mathrm{comp}}(\rho^{(k+1)})$ нельзя получить из старых энергий~\eqref{eq:simp-example-energies}: после изменения плотностей нужно собрать новую $K$, решить новую систему $KU=F$ и только затем вычислить податливость. Именно это отличает физическую оптимизацию от однократного ранжирования элементов по чувствительности.

Один шаг описанного алгоритма можно рассматривать как отображение
\begin{equation}
\label{eq:simp-refinement-map}
 \mathcal R_{\mathrm{SIMP}}
 :
 (\rho,c)
 \longmapsto
 \rho^+.
\end{equation}
После $K$ шагов
\begin{equation}
 \rho^{(K)}
 =
 \mathcal R_{\mathrm{SIMP}}^{(K)}
 \bigl(\rho^{(0)},c\bigr).
\end{equation}
Такое представление важно для дальнейшей главы: начальная плотность $\rho^{(0)}$ не обязана быть равномерной или случайной. Её может предложить генеративная модель.

Классический запуск и генеративный запуск тогда отличаются только инициализацией:
\begin{equation}
 \rho_{\mathrm{base}}^{(0)}
 \xmapsto{\ \mathcal R_{\mathrm{SIMP}}^{(K)}\ }
 \rho_{\mathrm{base}}^{(K)},
 \qquad
 G_\phi(z,c)
 \xmapsto{\ \mathcal R_{\mathrm{SIMP}}^{(K)}\ }
 \rho_{\mathrm{gen}}^{(K)}.
\end{equation}
Это даёт проверяемый смысл фразе ``генератор ускоряет проектирование''. Недостаточно получить визуально убедительный $G_\phi(z,c)$. Нужно сравнить стоимость достижения заданного физического качества после одинаковой независимой проверки: число дорогих решённых задач, итоговую податливость, соблюдение ограничений и устойчивость результата при изменении сетки.

Физический критерий и правила допустимости теперь зафиксированы. Генеративный этап меняет только механизм получения начальных кандидатов: вместо одной траектории оптимизации рассматривается условное распределение конструкций $p_\phi(\theta\mid c)$.

\section{Генеративная модель как распределение конструкций}

До этого момента конструкция возникала как одна точка пространства проектных переменных. Классический алгоритм получает начальное состояние $\theta^{(0)}$ и строит одну последовательность
\begin{equation}
 \theta^{(0)},\theta^{(1)},\ldots,\theta^{(K)}.
\end{equation}
Генеративная постановка меняет этот объект. При фиксированном инженерном условии $c$ требуется не единственная точка, а распределение допустимых или перспективных конструкций
\begin{equation}
\label{eq:conditional-design-distribution}
 \theta\sim p_\phi(\theta\mid c).
\end{equation}
Параметры $\phi$ принадлежат обучаемой модели, а $c$ описывает саму проектную задачу: нагрузку, закрепления, разрешённый объём материала и другие условия, которые должны сохранять физический смысл при генерации.

Это различие принципиально. Если две конструкции имеют близкую податливость, но разную геометрию, классическая локальная оптимизация может обнаружить только одну из них из выбранной инициализации. Модель~\eqref{eq:conditional-design-distribution} должна уметь порождать несколько вариантов при одном и том же $c$. Поэтому случайность здесь нужна не ради шума как такового, а ради исследования множества конкурирующих решений.

В этой секции diffusion-модель не выводится заново. Глава~\ref{chap:stochastic-differential-equations} уже дала прямой процесс зашумления и обратную динамику. Здесь рассматривается инженерный интерфейс: что считается состоянием, что входит в условие и каким образом физический критерий может менять распределение обратного шага.

Пусть $\mathcal D=\{(\theta^{(i)},c^{(i)})\}_{i=1}^{n}$ --- набор ранее полученных конструкций и соответствующих им проектных условий. В обычной регрессии модель могла бы пытаться восстановить одну функцию
\begin{equation}
\label{eq:single-design-regression}
 \widehat\theta=f_\phi(c).
\end{equation}
Такое отображение удобно, если каждому условию фактически соответствует один интересующий ответ. В генеративном проектировании это ограничение избыточно: для одного $c$ могут существовать геометрически разные конструкции с сопоставимым физическим качеством.

Поэтому вместо~\eqref{eq:single-design-regression} вводится случайная переменная
\begin{equation}
 z\sim p_Z,
 \qquad
 \theta=G_\phi(z,c),
\end{equation}
которая индуцирует условное распределение $p_\phi(\theta\mid c)$. При фиксированном $c$ разные значения $z$ должны давать разные кандидаты. При фиксированном $z$ изменение $c$ должно менять конструкцию в соответствии с новой постановкой задачи.

Это сразу задаёт два независимых требования. Во-первых, распределение должно зависеть от инженерного условия. Во-вторых, оно не должно вырождаться в одну и ту же геометрию для всех $z$. Первое относится к conditioning, второе --- к разнообразию. Ни одно из этих свойств само по себе не гарантирует низкую податливость или технологическую допустимость.

Работа Oh et al. нужна здесь как переход от обычной топологической оптимизации к исследованию множества проектов: целью становится получение нескольких вариантов, а не одного результата оптимизатора. У Mazé и Ahmed этот переход реализован через условную diffusion-модель, которая генерирует топологии под заданные ограничения.

\subsection{Что является условием инженерной генерации}

Запись $p_\phi(\theta\mid c)$ полезна только тогда, когда $c$ разложено на физически интерпретируемые компоненты. Для задачи статической топологической оптимизации удобно писать
\begin{equation}
 c=\bigl(c_{\mathrm{load}},c_{\mathrm{bc}},f_{\mathrm{vol}},c_{\mathrm{geom}},c_{\mathrm{mat}}\bigr),
\end{equation}
где $c_{\mathrm{load}}$ задаёт нагрузки, $c_{\mathrm{bc}}$ --- закрепления, $f_{\mathrm{vol}}$ --- допустимую долю материала, $c_{\mathrm{geom}}$ --- неизменяемую часть расчётной области, а $c_{\mathrm{mat}}$ --- фиксированные свойства материала, если они варьируются между задачами.

Не обязательно передавать все компоненты одной числовой строкой. Пространственное условие естественно представлять полем на той же сетке, что и конструкцию. Например, нагрузка в узле может кодироваться двумя каналами
\begin{equation}
 L_x(x_j)=F_{j,x},
 \qquad
 L_y(x_j)=F_{j,y},
\end{equation}
а фиксированная доля материала --- постоянным каналом
\begin{equation}
 V(x_j)=f_{\mathrm{vol}}
 \quad\text{для всех }x_j.
\end{equation}
Такой способ сохраняет пространственную привязку условия к расчётной области.

В датасете TopoDiff каждая обучающая запись содержит бинарную топологию, заданную долю материала, поле напряжения Мизеса, поле плотности энергии деформации и два канала нагрузки. Сами конструкции получены SIMP для разных комбинаций доли материала, закреплений и нагрузок. Это важная инженерная деталь: генератор учится не на абстрактных изображениях, а на результатах определённой физической постановки.

Отсюда следует правило для нашего контура. Нельзя менять кодировку $c$ между обучением и применением, сохраняя прежнюю интерпретацию модели. Если, например, величина силы при обучении нормировалась одним масштабом, а при инференсе другим, генератор получает другую задачу, даже если численно массивы выглядят похожими.

\subsection{Состояние диффузионной модели как конструкция}

Для плотностного представления состояние обратного процесса обозначим через
\begin{equation}
 x_t\in\mathbb R^N,
\end{equation}
где координаты соответствуют элементам или пикселям проектной области. При $t=T$ состояние близко к шуму, а при $t=0$ должно кодировать конструкцию. Для бинарного результата после генерации дополнительно может применяться порог или другое отображение к допустимому представлению.

Условный обратный переход записывается в общем виде как
\begin{equation}
\label{eq:conditional-reverse-transition}
 p_\phi(x_{t-1}\mid x_t,c).
\end{equation}
Если переход гауссовский,
\begin{equation}
\label{eq:gaussian-reverse-transition}
 p_\phi(x_{t-1}\mid x_t,c)
 =
 \mathcal N\!\left(
  \mu_\phi(x_t,c),
  \Sigma_t
 \right).
\end{equation}
Конструкция появляется не одним вызовом сети, а последовательностью условных переходов
\begin{equation}
 x_T\to x_{T-1}\to\cdots\to x_1\to x_0=\theta.
\end{equation}
Именно последовательная природа позволяет корректировать траекторию дополнительным guidance на промежуточных шагах.

При этом $x_t$ для больших $t$ ещё не является физически осмысленной конструкцией. На сильно зашумлённом состоянии нельзя буквально собирать FEM-модель и интерпретировать результат как реальную деталь. Если guidance использует суррогат физической величины, этот суррогат должен быть обучен работать и на промежуточных шумных состояниях. Mazé и Ahmed отдельно учитывают это требование для регрессора compliance.

\subsection{Условие генерации и направляющая поправка}

Conditioning означает, что условие $c$ входит непосредственно в обучаемый обратный переход~\eqref{eq:conditional-reverse-transition}. Модель учится порождать разные распределения для разных задач. Схематически
\begin{equation}
 c_1\neq c_2
 \quad\Longrightarrow\quad
 p_\phi(\theta\mid c_1)
 \neq
 p_\phi(\theta\mid c_2).
\end{equation}
Guidance решает другую задачу. Уже имеющийся переход дополнительно смещается в сторону конструкций, которые лучше по некоторому внешнему критерию.

Пусть базовый условный шаг имеет плотность $p_\phi(x_{t-1}\mid x_t,c)$, а внешняя функция $q_\psi(x_{t-1},c)$ штрафует нежелательное свойство. Тогда направленное распределение можно записать как
\begin{equation}
 \widetilde p(x_{t-1}\mid x_t,c)
 \propto
 p_\phi(x_{t-1}\mid x_t,c)
 \exp\!\bigl(-\alpha q_\psi(x_{t-1},c)\bigr),
\end{equation}
где $\alpha\geq0$ задаёт силу guidance. При $\alpha=0$ остаётся исходная модель. При большом $\alpha$ распределение сильнее концентрируется около малых значений $q_\psi$, но одновременно может потерять разнообразие или выйти в область, где сам суррогат ненадёжен.

В TopoDiff conditioning и guidance действительно разделены. Условная модель использует информацию о нагрузках, доле материала и физических полях, а дополнительное направление применяется при генерации для уменьшения compliance и floating material. Поэтому нельзя называть любой вход модели ``guidance'': часть информации задаёт саму условную задачу, а часть изменяет предпочтение внутри уже заданной задачи.

\subsection{Физические поля как часть условия}

Одни только координаты нагрузок и закреплений не показывают модели, как эти условия распространяются по области. TopoDiff дополнительно использует поля, рассчитанные для полной области под заданными нагрузками и граничными условиями. В частности, вводится напряжение Мизеса
\begin{equation}
 \sigma_{\mathrm{vm}}
 =
 \sqrt{
  \sigma_{11}^{2}
  -\sigma_{11}\sigma_{22}
  +\sigma_{22}^{2}
  +3\sigma_{12}^{2}
 },
\end{equation}
и плотность энергии деформации
\begin{equation}
 W
 =
 \frac12
 \left(
  \sigma_{11}\varepsilon_{11}
  +\sigma_{22}\varepsilon_{22}
  +2\sigma_{12}\varepsilon_{12}
 \right).
\end{equation}
Эти поля несут пространственную информацию о реакции расчётной области на условия задачи.

Преимущество такого conditioning понятно: вместо того чтобы самостоятельно вывести из нескольких точечных ограничений глобальную картину силового пути, модель получает уже вычисленные физические признаки. Но цена также очевидна: эти признаки требуют предварительного физического расчёта. Следовательно, ускорение генерации нельзя оценивать только временем обратного diffusion-процесса.

Если обозначить время подготовки условия через $T_{\mathrm{prep}}$, а собственно генерации через $T_{\mathrm{gen}}$, то честная стоимость одного нового сценария имеет вид
\begin{equation}
 T_{\mathrm{infer}}
 =
 T_{\mathrm{prep}}
 +T_{\mathrm{gen}}
 +T_{\mathrm{verify}}.
\end{equation}
Последний член появляется потому, что сгенерированная конструкция всё равно требует независимой физической проверки. Направление по суррогату используется только при генерации и не заменяет последующий расчёт и SIMP-уточнение.

\subsection{Направление по суррогатной податливости}

Mazé и Ahmed вводят отдельный регрессор $\widehat C_\psi$, который по промежуточной топологии и условию оценивает compliance. Направленный переход задаётся произведением базовой плотности на экспоненциальный физический вес
\begin{equation}
 \widetilde p_{\phi,\psi}(x_{t-1}\mid x_t,c)
 \propto
 p_\phi(x_{t-1}\mid x_t,c)
 \exp\!\left(-\alpha\widehat C_\psi(x_{t-1},c)\right).
\end{equation}
Чем выше предсказанная податливость, тем меньше относительный вес состояния.

Для гауссовского перехода~\eqref{eq:gaussian-reverse-transition} линейное разложение регрессора около текущего среднего даёт практически важную формулу. Обозначим
\begin{equation}
 g_C
 =
 \nabla_x\widehat C_\psi(x,c)
 \big|_{x=\mu_\phi}.
\end{equation}
Тогда новое среднее в первом приближении равно
\begin{equation}
\label{eq:guided-gaussian-mean}
 \boxed{
 \mu_{\mathrm{guided}}
 =
 \mu_\phi
 -\alpha\Sigma_t g_C
 }.
\end{equation}
Таким образом, guidance не заменяет diffusion-модель новым оптимизатором. Он локально сдвигает каждый стохастический шаг в направлении уменьшения предсказанной податливости.

Матрица $\Sigma_t$ здесь существенна. Если дисперсия по некоторой координате мала, одинаковый градиент вызывает малый сдвиг. Если большая --- сдвиг сильнее. Поэтому формула~\eqref{eq:guided-gaussian-mean} учитывает геометрию самого перехода, а не просто вычитает градиент фиксированного масштаба.

Наконец, $\widehat C_\psi$ остаётся суррогатом. Низкое значение регрессора не является доказательством низкой истинной податливости. После генерации требуется решить $K(\theta)U=F$ и вычислить физический $J_{\mathrm{comp}}$. В противном случае система оценивает сама себя той же приближённой моделью, которая направляла генерацию.

Податливость не обнаруживает все геометрические дефекты. В дискретной топологии могут появиться компоненты материала, не связанные с основной несущей конструкцией. В TopoDiff это свойство называется floating material и рассматривается как отдельный индикатор технологической недопустимости.

Пусть классификатор возвращает вероятность отсутствия такого дефекта
\begin{equation}
 P_\chi(y_{\mathrm{ok}}=1\mid x_t,c).
\end{equation}
Тогда направление к более желательному классу определяется градиентом
\begin{equation}
 g_{\mathrm{fm}}
 =
 \nabla_{x_t}
 \log P_\chi(y_{\mathrm{ok}}=1\mid x_t,c).
\end{equation}
Вместе с compliance-guidance обратный шаг получает два разных сигнала: один отвечает за функциональный критерий, другой --- за геометрическую связанность.

Это разделение полезно методологически. Нельзя заменить все инженерные требования одним словом ``качество''. Например, конструкция может иметь малую податливость, но содержать отдельный островок материала; другая может быть связной, но нарушать заданную долю материала. Эти ошибки требуют разных проверок и, при необходимости, разных направляющих функций.

Сила каждого guidance является гиперпараметром. В TopoDiff масштабы для compliance и floating material подбираются по отдельной validation-выборке. Это правильный протокол: значения нельзя выбирать по итоговому test-набору, иначе измеренная эффективность будет оптимистично смещена.

\subsection{Законченный расчёт: один направленный diffusion-шаг}

Рассмотрим двумерное промежуточное состояние. Пусть условная diffusion-модель без guidance предлагает переход
\begin{equation}
 x_{t-1}
 \sim
 \mathcal N(\mu,\Sigma),
 \qquad
 \mu=
 \begin{pmatrix}
 0.6\\
 0.4
 \end{pmatrix},
 \qquad
 \Sigma=0.04I.
\end{equation}
Пусть локальный surrogate compliance имеет простой вид
\begin{equation}
 \widehat C(x)
 =
 \frac12(x_1-1)^2
 +\frac12x_2^2.
\end{equation}
Его минимум расположен в точке $(1,0)^{\trans}$. В текущем среднем
\begin{equation}
 g_C
 =
 \nabla\widehat C(\mu)
 =
 \begin{pmatrix}
 -0.4\\
 0.4
 \end{pmatrix}.
\end{equation}
При $\alpha=1$ формула~\eqref{eq:guided-gaussian-mean} даёт
\begin{equation}
 \begin{aligned}
 \mu_{\mathrm{guided}}
 &=
 \begin{pmatrix}
 0.6\\
 0.4
 \end{pmatrix}
 -0.04
 \begin{pmatrix}
 -0.4\\
 0.4
 \end{pmatrix}\\
 &=
 \boxed{
 \begin{pmatrix}
 0.616\\
 0.384
 \end{pmatrix}}
 .
 \end{aligned}
\end{equation}
Среднее действительно сдвинулось к точке $(1,0)^{\trans}$.

До guidance surrogate даёт
\begin{equation}
 \widehat C(\mu)
 =
 \frac12(0.4^2+0.4^2)
 =0.16.
\end{equation}
После сдвига
\begin{equation}
 \begin{aligned}
 \widehat C(\mu_{\mathrm{guided}})
 &=
 \frac12(0.384^2+0.384^2)\\
 &=0.147456.
 \end{aligned}
\end{equation}
Локальная surrogate-оценка уменьшилась на
\begin{equation}
 \frac{0.16-0.147456}{0.16}
 \approx
 7.84\%.
\end{equation}

Этот расчёт показывает только механизм~\eqref{eq:guided-gaussian-mean}. Из него нельзя заключить, что истинный физический критерий улучшился на $7.84\%$: для этого двухкоординатный $x$ должен быть декодирован в конструкцию, затем требуется физический решатель. Аналогично нельзя заключить, что новый кандидат допустим по объёму или связанности. Guidance меняет вероятность траектории, но не отменяет независимые инженерные ограничения.

Генеративный блок можно описать интерфейсом
\begin{equation}
 \mathcal G_\phi:
 (c,z_1,\ldots,z_M)
 \longmapsto
 \{\theta^{(m)}\}_{m=1}^{M}.
\end{equation}
На входе находится одна инженерная задача и несколько реализаций случайности. На выходе --- набор кандидатов, а не окончательный ответ.

Для каждого кандидата ещё неизвестны истинное значение $J_h$, точное выполнение ограничений, устойчивость к уточнению сетки и преимущество относительно обычного SIMP-старта. Поэтому правильный следующий оператор имеет вид
\begin{equation}
 \{\theta^{(m)}\}_{m=1}^{M}
 \longrightarrow
 \text{физическая оценка}
 \longrightarrow
 \text{отбор}
 \longrightarrow
 \text{уточнение}.
\end{equation}

Условное распределение лишь создаёт кандидатов. Итоговый проект определяется последующим независимым физическим расчётом и классической оптимизацией, а не значением суррогата внутри генератора.

\section{Генерация, физика и оптимизация в одном контуре}
\label{sec:generative-optimization-loop}

Генеративная модель становится полезной для инженерного проектирования только после того, как её выход включён в тот же физический контур, по которому проверяется обычная конструкция. Поэтому итоговый объект этой секции --- не отдельный генератор и не отдельный SIMP, а композиция
\begin{equation}
 (c,z)
 \xmapsto{\ G_\phi\ }
 \theta^{(0)}
 \xmapsto{\ \mathcal P_h\ }
 \widetilde\theta^{(0)}
 \xmapsto{\ \mathcal R_{\mathrm{opt}}^{(K)}\ }
 \theta^{(K)}
 \xmapsto{\ S_{h/2}\ }
 \bigl(U_{h/2},J_{h/2}\bigr).
\end{equation}
Здесь $\mathcal P_h$ обозначает не обучение, а инженерную подготовку кандидата к точному расчёту: приведение к допустимому представлению, проверку обязательных геометрических условий и, если нужно, дешёвую предварительную оценку. Оператор $\mathcal R_{\mathrm{opt}}^{(K)}$ выполняет $K$ шагов физической оптимизации. Последний оператор $S_{h/2}$ не участвует в выборе кандидата и служит независимой проверкой.

Такая постановка снимает ложное противопоставление ``нейросеть или физика''. Генератор решает задачу поиска хорошей начальной области, а физический оптимизатор решает локальную задачу улучшения при явно заданных ограничениях. В работе Giannone и Ahmed эта идея реализована буквально: условная diffusion-модель выдаёт предварительную топологию, после чего несколько итераций SIMP уменьшают податливость и устраняют часть технологически плохих фрагментов. Авторы используют $5$ или $10$ итераций уточнения, а не полный повтор обычного длинного запуска SIMP.

\subsection{Генератор как инициализатор, а не замена решателя}

Пусть для фиксированного условия $c$ классический метод стартует из некоторой стандартной конструкции $\theta_{\mathrm{base}}^{(0)}$, например из равномерного поля плотности. Генеративный метод даёт другой старт
\begin{equation}
 \theta_{\mathrm{gen}}^{(0)}=G_\phi(z,c).
\end{equation}
После этого обе конструкции можно передать одному и тому же физическому оператору уточнения:
\begin{equation}
 \theta_{\mathrm{base}}^{(K)}
 =\mathcal R_{\mathrm{opt}}^{(K)}
 \bigl(\theta_{\mathrm{base}}^{(0)},c\bigr),
 \qquad
 \theta_{\mathrm{gen}}^{(K)}
 =\mathcal R_{\mathrm{opt}}^{(K)}
 \bigl(\theta_{\mathrm{gen}}^{(0)},c\bigr).
\end{equation}
В таком сравнении генеративная часть отвечает только за качество и разнообразие начальных точек. Это важно: если после генератора используется другой критерий, другая сетка или более сильный оптимизатор, улучшение уже нельзя приписывать одной генеративной модели.

Для topology optimization естественным физическим оператором является SIMP из раздела~\ref{sec:physics-based-design-optimization}. Тогда
\begin{equation}
\label{eq:generated-simp-start}
 \rho_{\mathrm{gen}}^{(0)}=G_\phi(z,c),
 \qquad
 \rho_{\mathrm{gen}}^{(K)}
 =\mathcal R_{\mathrm{SIMP}}^{(K)}
 \bigl(\rho_{\mathrm{gen}}^{(0)},c\bigr).
\end{equation}
Формула~\eqref{eq:generated-simp-start} даёт точный смысл словам ``генеративная оптимизация'': генератор не объявляет проект оптимальным, а меняет начальную точку детерминированного поиска.

Это особенно полезно в невыпуклой задаче. При SIMP-показателе $p>1$ приведённый критерий зависит от плотностей нелинейно, поэтому разные старты могут вести к разным локальным решениям. Генератор способен предложить несколько геометрически разных стартов под одно условие $c$; классическая оптимизация затем локально улучшает каждый из них по одному и тому же физическому критерию.

\subsection{Пул кандидатов и дешёвый предварительный отбор}

Главное преимущество условного генератора проявляется не в одном сэмпле, а в возможности получить несколько стартов для одной инженерной задачи. Поэтому при фиксированном $c$ генерируем $M$ независимых кандидатов с разными случайными кодами $z_m$:
\begin{equation}
 \mathcal C_M(c)
 =
 \left\{
  \theta^{(m)}=G_\phi(z_m,c)
 \right\}_{m=1}^{M},
 \qquad
 z_m\overset{\mathrm{iid}}{\sim}p_Z.
\end{equation}
Полный FEM-расчёт для всех $M$ кандидатов может уничтожить выигрыш по стоимости. Поэтому перед точной физикой допустим дешёвый фильтр, но его роль должна быть ограниченной: он сокращает число дорогих расчётов, а не заменяет итоговую проверку.

Для каждого кандидата вычислим набор дешёвых признаков
\begin{equation}
 q^{(m)}
 =
 \left(
  \widehat J_h^{(m)},
  \widehat r_{\mathrm{phys}}^{(m)},
  \eta^{(m)},
  e_{\mathrm{vol}}^{(m)}
 \right),
\end{equation}
где $\widehat J_h$ --- surrogate-оценка критерия, $\widehat r_{\mathrm{phys}}$ --- дешёвая оценка физической невязки, $\eta$ --- индикатор выхода из области применимости, а
\begin{equation}
 e_{\mathrm{vol}}(\theta;c)
 =
 \frac{\left|V(\theta)-V_{\mathrm{target}}(c)\right|}
 {V_{\mathrm{target}}(c)}
\end{equation}
измеряет нарушение ограничения на количество материала.

После этого выбирается подмножество
\begin{equation}
 \mathcal S_B(c)
 \subseteq
 \mathcal C_M(c),
 \qquad
 |\mathcal S_B(c)|=B\ll M,
\end{equation}
которое уже проходит точный физический расчёт. Значение $B$ является частью вычислительного бюджета. Если взять $B=M$, screening ничего не экономит. Если взять $B=1$, дешёвая модель получает слишком большую власть: одна ошибка ранжирования может удалить лучший кандидат до FEM.

Отбор должен сохранять не только малое $\widehat J_h$. Если несколько кандидатов почти одинаковы, нет смысла тратить дорогой физический расчёт на все копии. Поэтому часть бюджета можно отдать лучшим по суррогатному критерию, а часть --- геометрически отличающимся кандидатам; качество и разнообразие затем оцениваются раздельно.

\subsection{Точная физическая проверка кандидатов}

После screening каждый выбранный кандидат получает настоящий физический отклик
\begin{equation}
 R_h\bigl(U_h^{(m)},\theta^{(m)};c\bigr)=0,
 \qquad
 J_h^{(m)}
 =J_h\bigl(\theta^{(m)},U_h^{(m)};c\bigr).
\end{equation}
Именно здесь surrogate-оценка перестаёт быть основанием для инженерного вывода. Если
\begin{equation}
 \left|
  J_h^{(m)}-\widehat J_h^{(m)}
 \right|
 >\tau_J,
\end{equation}
то кандидат не ``почти хороший'', а плохо описанный быстрым фильтром. Его дальнейшая судьба определяется точной моделью.

Отдельно проверяются жёсткие условия. Обозначим через $I_{\mathrm{load}}(\theta;c)$ индикатор корректной передачи нагрузки на материал, а через $I_{\mathrm{geom}}(\theta;c)$ --- индикатор отсутствия запрещённых компонент. Тогда минимальную проверку можно записать как
\begin{equation}
 I_{\mathrm{hard}}(\theta;c)
 =
 \mathbf 1\!\left\{e_{\mathrm{vol}}\leq\tau_V\right\}
 I_{\mathrm{load}}(\theta;c)
 I_{\mathrm{geom}}(\theta;c).
\end{equation}
Формула не утверждает, что это исчерпывающий набор технологических ограничений. Она фиксирует принцип: каждое обязательное условие проверяется собственной вычислимой процедурой.

Направляющая поправка может уменьшать вероятность несвязного материала, но не гарантирует связность. Точно так же условие по доле материала не означает, что каждый отдельный образец соблюдает заданный объём. Точная проверка переводит статистическое предпочтение в решение: принять, отклонить или уточнить кандидат.

\subsection{Несколько шагов SIMP как физическое уточнение}

Если кандидат прошёл минимальную проверку, запускается короткое физическое уточнение. Для плотностной задачи один шаг уже определён формулой~\eqref{eq:simp-refinement-map}. Последовательность имеет вид
\begin{equation}
 \rho^{(0)}
 \xmapsto{\ \mathcal R_{\mathrm{SIMP}}\ }
 \rho^{(1)}
 \xmapsto{\ \mathcal R_{\mathrm{SIMP}}\ }
 \cdots
 \xmapsto{\ \mathcal R_{\mathrm{SIMP}}\ }
 \rho^{(K)},
\end{equation}
где на каждом шаге заново собирается $K(\rho^{(k)})$, решается равновесие и пересчитывается чувствительность. Поэтому уточнение действительно использует физику, а не косметически сглаживает изображение.

Giannone и Ahmed строят именно такой гибрид: TopoDiff-FF быстро генерирует предварительную топологию, затем $5$ или $10$ итераций SIMP используются для уменьшения compliance и исправления floating material. Смысл малого $K$ состоит в том, что генератор уже должен поставить систему в перспективную область. Если после каждого сэмпла требуется полный по длине обычный запуск оптимизатора, генеративный блок не доказал вычислительного преимущества.

Остановку удобно связывать не с фиксированным числом шагов, а также с физическим улучшением. Например,
\begin{equation}
 \frac{
  \left|J_h(\rho^{(k+1)})-J_h(\rho^{(k)})\right|
 }{
  \max\{1,|J_h(\rho^{(k)})|\}
 }
 <\varepsilon_J
\end{equation}
вместе с выполнением ограничений означает, что дальнейшее локальное уточнение даёт малый эффект. Для сравнения разных методов лучше фиксировать одинаковый бюджет либо одинаковый критерий остановки, а не выбирать более удобное условие отдельно для генеративного старта.

После уточнения кандидат получает две версии критерия:
\begin{equation}
 J_h^{(0)}=J_h\bigl(\rho^{(0)}\bigr),
 \qquad
 J_h^{(K)}=J_h\bigl(\rho^{(K)}\bigr).
\end{equation}
Их разность
\begin{equation}
 \Delta J_{\mathrm{ref}}
 =
 J_h^{(0)}-J_h^{(K)}
\end{equation}
показывает, какую часть качества дал локальный физический оптимизатор после генерации. Если $\Delta J_{\mathrm{ref}}$ систематически огромна, нельзя писать, что генератор сам выдаёт готовые физически качественные решения.

\subsection{Многостартовый поиск и сохранение разнообразия}

Пусть после точной проверки осталось $B$ стартов. Каждый уточняется независимо:
\begin{equation}
 \rho_m^{(K)}
 =
 \mathcal R_{\mathrm{SIMP}}^{(K)}
 \bigl(\rho_m^{(0)},c\bigr),
 \qquad
 m=1,\ldots,B.
\end{equation}
Если требуется только один проект, естественный итоговый выбор ---
\begin{equation}
 m^\star
 =
 \arg\min_{
  m:\,\rho_m^{(K)}\in\Theta_{\mathrm{adm}}(c)
 }
 J_{h/2}\bigl(\rho_m^{(K)};c\bigr).
\end{equation}
Но generative design часто нужен именно потому, что инженер хочет несколько вариантов. Тогда выбор одного $m^\star$ слишком рано уничтожает информацию о пространстве решений.

Поэтому полезно хранить множество
\begin{equation}
 \mathcal V(c)
 =
 \left\{
  \rho_m^{(K)}:
  \rho_m^{(K)}\in\Theta_{\mathrm{adm}}(c),
  \ J_{h/2}(\rho_m^{(K)})\leq J_{\max}
 \right\}.
\end{equation}
Внутри $\mathcal V(c)$ конструкции уже удовлетворяют физическому порогу, поэтому разнообразие можно обсуждать как дополнительное свойство, а не как замену качеству.

Это и есть отличие многостартовой генеративной схемы от одного запуска SIMP. Генератор полезен не тем, что рисует одну структуру быстрее, а тем, что может дать несколько разных областей притяжения. Короткое уточнение затем переводит эти предложения в пространство физически проверенных решений. Работа Regenwetter et al. позднее позволит измерить, действительно ли полученное множество отличается по novelty и diversity, а не только визуально кажется разнообразным.

\subsection{Вычислительный бюджет и условие полезности гибридной схемы}

Сравнивать методы только по времени одного вызова генератора неверно. Полная стоимость должна включать всё, что требуется для получения проверенного проекта. Для одного нового условия $c$ запишем
\begin{equation}
 C_{\mathrm{hyb}}(c)
 =
 C_{\mathrm{cond}}
 +M C_{\mathrm{gen}}
 +M C_{\mathrm{screen}}
 +B C_{\mathrm{FEM}}
 +B K C_{\mathrm{iter}}
 +B C_{\mathrm{verify}}.
\end{equation}
Здесь $C_{\mathrm{cond}}$ --- стоимость подготовки conditioning, $C_{\mathrm{gen}}$ --- одного сэмпла, $C_{\mathrm{screen}}$ --- дешёвого отбора, $C_{\mathrm{FEM}}$ --- первого точного расчёта, $C_{\mathrm{iter}}$ --- одного шага уточнения, а $C_{\mathrm{verify}}$ --- независимой финальной проверки.

Для обычного многостартового оптимизатора аналогичная стоимость равна
\begin{equation}
 C_{\mathrm{base}}(c)
 =
 B_0 K_0 C_{\mathrm{iter}}
 +B_0 C_{\mathrm{verify}},
\end{equation}
если conditioning и генерация отсутствуют. Гибридная схема имеет практический смысл, когда при сопоставимом итоговом качестве
\begin{equation}
 C_{\mathrm{hyb}}(c)<C_{\mathrm{base}}(c)
\end{equation}
или когда при сопоставимой стоимости она даёт лучшее множество проверенных решений.

К стоимости эксплуатации нельзя бесследно спрятать обучение. Если генератор обучался один раз за цену $C_{\mathrm{train}}$ и затем применяется к $N_{\mathrm{use}}$ новым задачам, амортизированная стоимость на задачу составляет
\begin{equation}
 \overline C_{\mathrm{hyb}}
 =
 \frac{C_{\mathrm{train}}}{N_{\mathrm{use}}}
 +\frac1{N_{\mathrm{use}}}
 \sum_{i=1}^{N_{\mathrm{use}}}
 C_{\mathrm{hyb}}(c_i).
\end{equation}
Поэтому генеративный подход особенно естественен для семейства повторяющихся задач. Для единственного расчёта дорогая подготовка датасета и обучение могут быть бессмысленнее прямой оптимизации.

Работа Giannone и Ahmed показывает ещё одну важную деталь стоимости. Дорогими могут быть не только FEM-итерации после генерации, но и подготовка условия модели. В TopoDiff-FF физические поля, для которых нужен предварительный FEA, заменяются дешёвой kernel-relaxation. Поэтому ускорение нужно считать от входного условия $c$ до проверенного проекта, а не от заранее подготовленного тензора до изображения.

\subsection{Выход из распределения и безопасный возврат к базовому решателю}

Генератор обучается на конечном наборе условий. Пусть $\eta(c)$ --- индикатор того, насколько новая постановка похожа на область обучения. Тогда гибридный метод должен иметь явную ветвь отказа:
\begin{equation}
 \eta(c)>\tau_{\mathrm{ood}}
 \quad\Longrightarrow\quad
 \theta^{(0)}=\theta_{\mathrm{base}}^{(0)}.
\end{equation}
Это не означает, что OOD-кандидат обязательно плох. Это означает только отсутствие основания использовать обученную модель как ускоритель без дополнительной проверки.

Возможна и более мягкая схема. Генеративный старт сохраняется как один из нескольких стартов, но к нему добавляется стандартная инициализация:
\begin{equation}
 \mathcal I(c)
 =
 \left\{
  \rho_{\mathrm{base}}^{(0)},
  G_\phi(z_1,c),\ldots,G_\phi(z_M,c)
 \right\}.
\end{equation}
Тогда генератор не может сделать алгоритм хуже только из-за того, что вытеснил базовую ветвь; финальный выбор по точной физике может вернуть обычное решение.

Именно OOD-режим показывает, зачем после генерации нужен оптимизатор. В экспериментах Giannone и Ahmed вариант TopoDiff-FF без уточнения заметно теряет качество при новых сочетаниях ограничений, тогда как добавление десяти шагов SIMP резко уменьшает compliance error и улучшает выполнение ограничений. Это эмпирический результат конкретной статьи, а не универсальная гарантия: в другой задаче локальный оптимизатор может не исправить плохую инициализацию за малое число шагов.

\subsection{Законченный расчёт: цена десяти шагов физического уточнения}

Рассмотрим опубликованный авторами режим с out-of-distribution ограничениями, обозначенный как task-2. Для TopoDiff-FF без SIMP средний compliance error равен
\begin{equation}
 \mathrm{CE}_0=58.36\%,
\end{equation}
а полное время inference в таблице составляет
\begin{equation}
 T_0=2.35\ \mathrm{s}.
\end{equation}
После десяти итераций SIMP авторы получают
\begin{equation}
 \mathrm{CE}_{10}=7.84\%,
 \qquad
 T_{10}=2.55\ \mathrm{s}.
\end{equation}
Сначала посчитаем относительное уменьшение ошибки податливости:
\begin{equation}
 \frac{\mathrm{CE}_0-\mathrm{CE}_{10}}{\mathrm{CE}_0}
 =
 \frac{58.36-7.84}{58.36}
 \approx
 \boxed{0.866}.
\end{equation}
То есть в этом эксперименте десять шагов физического уточнения уменьшают средний compliance error примерно на $86.6\%$ относительно варианта без уточнения.

Дополнительное время равно
\begin{equation}
 \Delta T
 =2.55-2.35
 =0.20\ \mathrm{s}.
\end{equation}
Относительное увеличение времени составляет
\begin{equation}
 \frac{\Delta T}{T_0}
 =
 \frac{0.20}{2.35}
 \approx
 \boxed{8.51\%}.
\end{equation}
Одновременно в той же таблице volume fraction error уменьшается с $1.97\%$ до $1.29\%$, а доля floating material --- с $7.86\%$ до $6.53\%$.

Этот расчёт не доказывает, что ``десять шагов SIMP всегда выгодны''. Он показывает более узкий и полезный факт: в конкретном OOD-эксперименте слабый генеративный кандидат был хорошей точкой для короткой физической коррекции. Поэтому сравнивать только TopoDiff-FF и обычный SIMP недостаточно; отдельным методом является именно композиция
\begin{equation}
 \boxed{
  G_\phi
  \;\longrightarrow\;
  \text{точная физика}
  \;\longrightarrow\;
  \mathcal R_{\mathrm{SIMP}}^{(K)}
  \;\longrightarrow\;
  \text{независимая проверка}
 }.
\end{equation}

Рабочий контур теперь полностью определён: генератор создаёт пул кандидатов, предварительный отбор сокращает число дорогих проверок, точный решатель отбраковывает физически несостоятельные варианты, короткий оптимизационный цикл улучшает прошедшие кандидаты, а OOD-ветвь сохраняет возможность вернуться к базовому методу. Осталось задать независимую процедуру оценки, в которой качество, допустимость, разнообразие и вычислительная стоимость измеряются отдельно.

\section{Проверка генеративной системы и полный проектный протокол}

Кандидат, прошедший точный физический расчёт и короткое оптимизационное уточнение, ещё не характеризует метод в целом. Генеративный алгоритм выдаёт множество конструкций: оно может быть похоже на обучающие данные, но нарушать ограничения; давать хорошие значения податливости, но почти не создавать новых геометрий; быть разнообразным, но игнорировать заданную долю материала. Эти свойства проверяются раздельно.

Regenwetter et al. предлагают именно такое разбиение. Для инженерных генеративных моделей они различают statistical similarity, design exploration, constraint satisfaction, design quality и conditioning. Для нашей задачи это означает: отдельно проверять похожесть на известные конструкции, новизну и разнообразие, физическую допустимость, значение целевого функционала и соответствие заданному условию $c$. Ни одна из этих координат не заменяет остальные.

\subsection{Что именно является объектом проверки}

Пусть тестовый набор содержит $N_{\mathrm{test}}$ инженерных постановок
\begin{equation}
 \mathcal C_{\mathrm{test}}
 =\{c_i\}_{i=1}^{N_{\mathrm{test}}}.
\end{equation}
Для каждой постановки генерируем $M$ независимых кандидатов
\begin{equation}
 \mathcal G_i
 =\left\{
  \theta_{ij}=G_\phi(z_{ij},c_i)
 \right\}_{j=1}^{M},
 \qquad
 z_{ij}\overset{\mathrm{iid}}{\sim}p_Z.
\end{equation}
После обязательной физической проверки и, если это часть исследуемого метода, фиксированного уточнения получаем множество
\begin{equation}
 \mathcal V_i
 =\left\{
  \widetilde\theta_{ij}
 \right\}_{j=1}^{M}.
\end{equation}
Важно не смешивать $\mathcal G_i$ и $\mathcal V_i$. Первая совокупность характеризует непосредственный выход генератора, вторая --- качество всей системы
\begin{equation}
 \mathcal A
 =
 \text{generator}
 \circ
 \text{physics gate}
 \circ
 \text{refinement}.
\end{equation}
Если в статье заявляется преимущество полного generative-design pipeline, основные инженерные метрики должны считаться на $\mathcal V_i$. Если требуется понять вклад самого генератора, дополнительно сравниваются показатели до и после уточнения.

Такое различие совпадает с разделением point- и set-метрик у Regenwetter et al. Point-метрика относится к отдельной конструкции, например её compliance или нарушение объёма. Set-метрика характеризует совокупность, например diversity. Среднее point-метрик и set-метрика --- не одно и то же: усреднение индивидуального качества не сообщает, насколько разнообразны решения.

\subsection{Допустимость и физическая состоятельность}

Для каждого проверенного кандидата определим индикатор выполнения жёстких условий
\begin{equation}
 I_{\mathrm{feas}}(\theta;c)
 =
 \mathbf 1\!\left\{
  \theta\in\Theta_{\mathrm{adm}}(c)
 \right\}.
\end{equation}
Тогда доля допустимых конструкций в тестовом эксперименте равна
\begin{equation}
 r_{\mathrm{feas}}
 =
 \frac{1}{N_{\mathrm{test}}M}
 \sum_{i=1}^{N_{\mathrm{test}}}
 \sum_{j=1}^{M}
 I_{\mathrm{feas}}(\widetilde\theta_{ij};c_i).
\end{equation}
Эта величина отвечает на простой вопрос: как часто система вообще выдаёт конструкцию, которую разрешено рассматривать дальше.

Для topology optimization бинарного индикатора недостаточно. Две конструкции могут обе нарушать связность, но одна иметь один изолированный пиксель, а другая --- большую плавающую компоненту. Поэтому рядом с бинарной проверкой полезно хранить величину нарушения. Для доли плавающего материала обозначим
\begin{equation}
 q_{\mathrm{float}}(\theta)
 =
 \frac{N_{\mathrm{disconnected}}(\theta)}
 {N_{\mathrm{material}}(\theta)}.
\end{equation}
Граница допустимости соответствует $q_{\mathrm{float}}=0$. Regenwetter et al. подчёркивают, что одной бинарной отметки здесь мало: количественная доля несвязного материала показывает не только факт нарушения, но и его величину.

Аналогично нагрузочная точка должна быть соединена с несущим материалом. Обозначим
\begin{equation}
 I_{\mathrm{load}}(\theta;c)
 =
 \mathbf 1\!\left\{
  \text{нагрузка передаётся на связную материальную компоненту}
 \right\}.
\end{equation}
Вместо визуального осмотра в эксперименте должна быть зафиксирована конкретная вычислительная процедура построения компонент связности. Иначе одинаковая картинка может получить разные оценки в двух реализациях.

Физическая состоятельность проверяется уже не геометрическим тестом, а решателем. Для каждой прошедшей конструкции решается
\begin{equation}
 R_{h_{\mathrm{val}}}
 \bigl(U_{h_{\mathrm{val}}},\theta;c\bigr)=0,
\end{equation}
причём $h_{\mathrm{val}}$ и параметры решателя фиксируются до сравнения методов. Если генератор обучался по данным на сетке $h$, разумно хотя бы для финального набора использовать независимую или более точную дискретизацию $h_{\mathrm{val}}$. Тогда маленькая ошибка на обучающем solver-е не принимается автоматически за физическое качество.

\subsection{Качество конструкции относительно физического базиса}

Для минимизируемого критерия введём базовое решение $\theta_{\mathrm{base}}(c)$, полученное заранее выбранным детерминированным методом. В topology optimization таким базисом естественно служит SIMP при фиксированных настройках. Относительный разрыв по критерию определим как
\begin{equation}
 \Delta_J(\theta;c)
 =
 \frac{
  J_{h_{\mathrm{val}}}(\theta;c)
  -
  J_{h_{\mathrm{val}}}(\theta_{\mathrm{base}};c)
 }{
  \left|J_{h_{\mathrm{val}}}(\theta_{\mathrm{base}};c)\right|
 }.
\end{equation}
Знак здесь содержателен. Значение $\Delta_J>0$ означает, что кандидат хуже базиса; $\Delta_J<0$ --- что он превзошёл выбранное базовое решение на той же проверочной физике.

В работах по генеративной topology optimization часто используется compliance error относительно SIMP-решения. Regenwetter et al. интерпретируют его как нормированное расстояние до целевого значения в пространстве физических характеристик. Для нашей записи
\begin{equation}
 e_C(\theta;c)
 =
 \frac{
  C(\theta;c)-C_{\mathrm{SIMP}}(c)
 }{
  C_{\mathrm{SIMP}}(c)
 }.
\end{equation}
Среднее $e_C$ показывает систематический сдвиг, медиана устойчивее к отдельным тяжёлым выбросам, а квантили показывают хвост плохих генераций. Поэтому один средний показатель недостаточен, если из тысячи кандидатов несколько десятков оказываются катастрофически мягкими.

Для полного гибридного метода дополнительно важен выигрыш от refinement. Для каждого кандидата запишем
\begin{equation}
 \Delta J_{\mathrm{ref},ij}
 =
 J_{h_{\mathrm{val}}}(\theta_{ij};c_i)
 -
 J_{h_{\mathrm{val}}}(\widetilde\theta_{ij};c_i).
\end{equation}
Если итоговый результат хорош, но почти весь выигрыш сосредоточен в $\Delta J_{\mathrm{ref}}$, вывод должен звучать как ``генератор дал полезную инициализацию для физического оптимизатора'', а не ``генератор решил задачу topology optimization''.

\subsection{Соответствие конструкции заданному условию}

Условная модель должна не просто выдавать допустимые конструкции, а реагировать на конкретное $c_i$. Пусть одна из компонент условия --- целевая доля материала $f_{\mathrm{vol}}(c)$. Для кандидата вычислим фактическую долю материала $\widehat f_{\mathrm{vol}}(\theta)$. Тогда conditioning adherence для этой компоненты можно записать как
\begin{equation}
 e_{\mathrm{vol}}(\theta;c)
 =
 \frac{
  \left|
   \widehat f_{\mathrm{vol}}(\theta)
   -f_{\mathrm{vol}}(c)
  \right|
 }{
  f_{\mathrm{vol}}(c)
 }.
\end{equation}
Regenwetter et al. прямо относят volume fraction error в topology generation к conditioning adherence, потому что volume fraction подаётся модели как условие, а затем может быть независимо вычислена по сгенерированной конструкции.

Для нагрузки и закреплений идея та же, но метрика зависит от представления. Если $c$ содержит координату приложения силы, нельзя считать выполнение условия по сходству входного тензора. Нужно проверить, что полученная конструкция физически соединяет область нагрузки с закреплением. Поэтому conditioning и constraint satisfaction могут использовать одну геометрическую процедуру, но отвечают на разные вопросы: первое проверяет соответствие заданному $c$, второе --- допустимость конструкции вообще.

Полезна проверка чувствительности к условию. Берём один и тот же латентный код $z$ и два условия $c_a\neq c_b$:
\begin{equation}
 \theta_a=G_\phi(z,c_a),
 \qquad
 \theta_b=G_\phi(z,c_b).
\end{equation}
Если изменение нагрузки, закрепления или объёма почти не меняет выход, модель может фактически игнорировать conditioning. Проверять это только по training loss нельзя: итоговая конструкция должна показать изменение измеримого инженерного свойства.

\subsection{Новизна и разнообразие после физического фильтра}

Generative design отличается от ускоренного предсказания одного SIMP-решения тем, что должен исследовать несколько областей пространства конструкций. Пусть $d(\theta_a,\theta_b)$ --- заранее определённое расстояние между представлениями конструкций. Для фиксированного условия $c_i$ простая средняя попарная дистанция равна
\begin{equation}
 D_i
 =
 \frac{2}{M(M-1)}
 \sum_{1\leq j<k\leq M}
 d(\widetilde\theta_{ij},\widetilde\theta_{ik}).
\end{equation}
Это мера разнообразия внутри одной генерации. Она не говорит, новы ли конструкции относительно известных решений.

Новизну относительно базиса можно измерить отдельно:
\begin{equation}
 N_{ij}
 =
 d\!\left(
  \widetilde\theta_{ij},
  \theta_{\mathrm{base}}(c_i)
 \right).
\end{equation}
В topology generation Regenwetter et al. предлагают именно distance to SIMP как показатель того, насколько генератор повторяет SIMP-решение или предлагает другой старт для повторной оптимизации.

Но большие $D_i$ и $N_{ij}$ полезны только после физического фильтра. Генератор легко может увеличить ``разнообразие'', выдавая повреждённые или несвязные конструкции. Поэтому для инженерного вывода рассматриваем
\begin{equation}
 \mathcal V_i^{\mathrm{good}}
 =
 \left\{
  \widetilde\theta_{ij}:
  I_{\mathrm{feas}}=1,
  \ \Delta_J\leq\tau_J
 \right\},
\end{equation}
и разнообразие считаем уже внутри $\mathcal V_i^{\mathrm{good}}$. Тогда ``разные'' означает разные среди физически конкурентных решений, а не среди любого выхода сети.

Такая последовательность не запрещает similarity-метрики. Они полезны для диагностики: слишком большое расстояние от обучающего распределения может означать сломанное представление. Но similarity остаётся вспомогательной координатой. На рисунке с велосипедными рамами у Regenwetter et al. две визуально очень близкие конструкции отличаются по прогибу более чем на два порядка, что как раз показывает, почему геометрическая похожесть не заменяет физическую оценку.

\subsection{Разделение сценариев и защита от утечки}

Случайный train/test split отдельных топологий может дать слишком лёгкую задачу. Если одна и та же комбинация нагрузки, закрепления и доли материала встречается в обеих частях, модель проверяется в основном на воспроизведение знакомого сценария. Для инженерного генератора важнее разделять не только изображения, но и постановки $c$.

Пусть множество всех условий разбито как
\begin{equation}
 \mathcal C
 =
 \mathcal C_{\mathrm{train}}
 \mathbin{\dot\cup}
 \mathcal C_{\mathrm{val}}
 \mathbin{\dot\cup}
 \mathcal C_{\mathrm{test}},
\end{equation}
причём финальные $c\in\mathcal C_{\mathrm{test}}$ не используются ни для обучения, ни для выбора гиперпараметров. Дополнительно полезно выделить OOD-набор $\mathcal C_{\mathrm{ood}}$, где меняется заранее определённый аспект постановки: например, положение нагрузки или комбинация закреплений выходит за обучающий диапазон.

Физический baseline для test и OOD также вычисляется без подглядывания генератора. Нельзя выбирать число SIMP-итераций или tolerances отдельно для каждого метода после просмотра результатов. До эксперимента фиксируются дискретизация, критерий остановки, допустимые ограничения, число сэмплов $M$, refinement-бюджет $K$ и вычислительный бюджет.

Отдельная опасность возникает, если surrogate использовался и для guidance, и для финальной оценки. Тогда система фактически проверяется собственным критиком. Поэтому финальное $J_{h_{\mathrm{val}}}$ вычисляется независимым физическим solver-ом. Тот же принцип уже встречался в главе~\ref{sec:generative-optimization-loop}: дешёвый screening может экономить вычисления, но не получает право выдавать окончательный инженерный вердикт.

\subsection{Законченный расчёт: почему одного рейтинга методов недостаточно}

Regenwetter et al. повторно оценивают два варианта TopoDiff. Первый получает условия из полей, рассчитанных методом конечных элементов, второй --- из дешёвого приближения. Для первого варианта медианная ошибка по податливости равна $1.99\%$, среднее расстояние до SIMP --- $11.6\%$, а среднее время генерации --- $5.87\ \mathrm{s}$. Для второго соответствующие величины равны $9.61\%$, $17.3\%$ и $2.35\ \mathrm{s}$.

Сначала сравним время. Относительное сокращение составляет
\begin{equation}
 \frac{5.87-2.35}{5.87}
 \approx
 \boxed{59.97\%}.
\end{equation}
По скорости kernel-conditioning заметно лучше.

Теперь новизна относительно SIMP:
\begin{equation}
 \frac{17.3-11.6}{11.6}
 \approx
 \boxed{49.14\%}.
\end{equation}
По этой метрике kernel-вариант также сильнее: его конструкции дальше от соответствующих SIMP-решений.

Но физическое качество меняется в противоположную сторону. Compliance error возрастает с $1.99\%$ до $9.61\%$:
\begin{equation}
 \frac{9.61-1.99}{1.99}
 \approx
 \boxed{382.9\%}.
\end{equation}
Следовательно, фраза ``kernel-вариант лучше FEA-варианта'' не следует из таблицы. Первый быстрее и новее относительно SIMP, второй точнее по compliance и лучше по ряду ограничений. Выбор зависит от роли генератора. Если нужен почти готовый проект, малый compliance error важнее. Если нужен дешёвый необычный старт для последующего SIMP, novelty и inference time могут иметь больший вес.

Этот расчёт показывает, почему нельзя сворачивать проверку в один произвольный score. Для инженерной системы естественнее хранить вектор
\begin{equation}
 \mathcal M
 =
 \left(
  r_{\mathrm{feas}},
  \operatorname{median} e_C,
  \overline e_{\mathrm{vol}},
  \overline N,
  \overline D,
  C_{\mathrm{online}}
 \right),
\end{equation}
а затем явно указывать, по каким координатам один метод выигрывает у другого.

После всех предыдущих секций проектный цикл можно записать одной композицией
\begin{equation}
 \boxed{
 c
 \longrightarrow
 G_\phi(z,c)
 \longrightarrow
 \text{screen}
 \longrightarrow
 S_h
 \longrightarrow
 \mathcal R_{\mathrm{opt}}^{(K)}
 \longrightarrow
 S_{h_{\mathrm{val}}}
 \longrightarrow
 \mathcal M
 }.
\end{equation}
Каждая стрелка имеет отдельный смысл. $G_\phi$ исследует пространство конструкций. Screening сокращает дорогой пул, но не заменяет физику. $S_h$ даёт точный для рабочего контура физический отклик. $\mathcal R_{\mathrm{opt}}^{(K)}$ локально улучшает кандидата. $S_{h_{\mathrm{val}}}$ выполняет независимую проверку. Вектор $\mathcal M$ сохраняет качество, допустимость, conditioning, разнообразие и стоимость раздельно.

Для воспроизводимого эксперимента вместе с результатом должны быть зафиксированы $\mathcal C_{\mathrm{train}}$, $\mathcal C_{\mathrm{val}}$, $\mathcal C_{\mathrm{test}}$, правило построения $\mathcal C_{\mathrm{ood}}$, число генераций $M$, случайные seed, физический solver и его сетка, baseline, число шагов уточнения $K$, определения всех ограничений и параметры метрик. Это соответствует замечанию Regenwetter et al.: настройки evaluation metrics и training data должны быть прозрачны, иначе сравнение моделей нельзя воспроизвести.

Для задач уровня AIRI содержательная связь теперь полностью определена без утверждений о конкретных внутренних системах. Машинное обучение отвечает за представление и быстрый поиск в большом пространстве кандидатов; механика и PDE задают проверяемые свойства; автоматическое дифференцирование и оптимизация дают путь уточнения; независимая валидация отделяет красивый или статистически правдоподобный объект от инженерно пригодной конструкции.

Глава 9 тем самым заканчивает теоретический путь от проектной переменной к проверенному множеству решений. Следующая глава не должна вводить ещё один слой общей теории. Её задача --- взять этот контракт и реализовать несколько законченных проектов, где видны данные, код, физический baseline, обучение, генерация, уточнение, метрики и итоговая проверка.
